\documentclass[11pt,a4paper]{article}

\usepackage[utf8]{inputenc}
\usepackage[T1]{fontenc}
\usepackage{lmodern}
\usepackage[a4paper,margin=1in]{geometry}
\usepackage{setspace}
\usepackage{amsmath,amssymb,bm}
\usepackage{booktabs}
\usepackage{threeparttable}
\usepackage{graphicx}
\usepackage{placeins}      % for \FloatBarrier
\usepackage{multirow}
\usepackage{array}
\usepackage{makecell}
\usepackage{caption}
\captionsetup[table]{skip=4pt}
\captionsetup[figure]{skip=6pt}

\usepackage{hyperref}
\usepackage[authoryear]{natbib}

\hypersetup{
    colorlinks=true,
    citecolor=blue,
    linkcolor=blue,
    urlcolor=blue
}
\onehalfspacing

\title{Asymmetric Volatility, Tail Risk, and Limits to Arbitrage in the Cross-Section of Stock Returns}
\author{Jasper Wolffenbuttel \and Koen Nijlaan \and Quinten van Amerongen \and Jules de Kruijk}
\date{\today}

\begin{document}

\maketitle

\begin{abstract}
We examine the mechanism driving the cross-sectional return predictability of realized
asymmetric volatility and introduce high-frequency realized expected shortfall (RES) as a
complementary intraday tail risk measure. Decomposing the Relative Signed Jump (RSJ) of
\citet{bollerslev2020} and RES into systematic and idiosyncratic components over a weekly
U.S.\ equity panel spanning 1993--2024, we find that virtually all predictive content
resides in the idiosyncratic channel. In quintile portfolio sorts, idiosyncratic RSJ earns
17--20 basis points per week --- comparable to the total RSJ spread --- while systematic RSJ
is economically negligible under both decomposition methods. Fama--MacBeth regressions
confirm this pattern: idiosyncratic RSJ commands average slopes of $-40$ to $-59$ basis
points per week, and the full decomposition raises the adjusted $R^2$ to 8.4\%,
versus 6.7\% for the aggregate specification. Formal Wald tests reject the joint null
that idiosyncratic components are unpriced but fail to reject the corresponding null for
systematic components. Crisis-state regressions reveal no amplification of the systematic
premia during NBER recessions, directly contradicting a rational risk-compensation
interpretation. By contrast, the idiosyncratic RSJ premium is significantly stronger
among illiquid, hard-to-arbitrage stocks, consistent with limits-to-arbitrage-induced
mispricing. Idiosyncratic RES independently captures tail severity beyond RSJ and is
itself tightly linked to trading frictions. Collectively, the evidence indicates that
firm-specific downside shocks generate return predictability that persists because
the stocks most exposed to idiosyncratic tail risk are precisely those where arbitrage
is most costly.
\end{abstract}

\newpage
\section{Introduction}
 
\citet{bollerslev2020} document that stocks whose intraday realized volatility is disproportionately driven by negative return innovations earn significantly higher future returns. This finding, based on the Relative Signed Jump (RSJ) measure, is striking: it implies that a quantity constructed entirely from high-frequency within-day variation contains economically meaningful information about the cross-section of weekly expected returns. Yet the mechanism behind this predictability remains unresolved. Does RSJ earn a risk premium because it captures exposure to adverse aggregate states of the world, or does it reflect firm-specific mispricing that persists because arbitrage is risky and costly?
 
This distinction has broad implications. If asymmetric volatility is priced because it reflects systematic downside exposure, then the RSJ effect represents compensation for a source of risk that investors cannot diversify away, and it should be most pronounced when aggregate economic conditions deteriorate. If, instead, the effect is driven by idiosyncratic shocks that are difficult to arbitrage, then it reflects persistent mispricing concentrated among stocks with high trading frictions \citep{shleifer1997,amihud2002}, and it should weaken or disappear among liquid, easily arbitraged stocks. Despite the importance of this question, the existing literature does not provide the tools to answer it. Two limitations stand in the way.
 
First, RSJ captures the directional asymmetry of volatility but is silent about the severity of extreme negative outcomes. A stock with many moderately negative intraday returns will exhibit a strongly negative RSJ, even if it never experiences deep tail losses. Conversely, a stock that suffers a few extreme intraday crashes may show a less negative RSJ if positive returns offset the asymmetry, despite facing substantial left-tail risk. Whether the return predictability associated with RSJ reflects exposure to extreme downside events or a broader pattern of volatility asymmetry therefore remains unclear. To address this, we introduce realized expected shortfall (RES), constructed from high-frequency returns under the self-similar scaling framework of \citet{Dimitriadis03072022}, which directly measures the average severity of intraday left-tail losses. If RES predicts future returns beyond RSJ, this would indicate that tail severity captures a dimension of downside risk that volatility asymmetry alone does not reflect.
 
Second, the aggregate RSJ measure confounds systematic and idiosyncratic sources of asymmetric volatility. A stock may exhibit negative RSJ either because its returns covary with market-wide downside movements, for instance during aggregate volatility spikes that disproportionately affect downside returns, or because it experiences firm-specific negative shocks unrelated to the broader market. These two sources have fundamentally different asset pricing implications: systematic asymmetric volatility should covary with aggregate state variables and command a risk premium, while idiosyncratic asymmetric volatility should not be priced under standard theory unless limits to arbitrage prevent its correction \citep{atilgan2020}. We decompose both RSJ and RES into systematic and idiosyncratic components using an intraday market model. For RSJ, we implement two alternative decomposition methods, one that separates intraday returns before constructing RSJ and one that decomposes the final RSJ series using a market-wide RSJ factor, to assess whether the source of predictability is sensitive to the decomposition approach.
 
These choices yield two testable implications that help identify the relative importance of each channel. Under a risk-based interpretation, the systematic components of RSJ and RES should drive predictability, with effects that intensify during NBER recession periods when exposure to aggregate downside risk matters most for asset prices. Under a mispricing-based interpretation, the idiosyncratic components should drive predictability, with effects that concentrate among illiquid or hard-to-arbitrage stocks where firm-specific shocks are most likely to persist. In practice, both channels may operate simultaneously; the empirical design is aimed at quantifying their relative contributions rather than ruling out either explanation entirely.
 
The central research question is:
 
\begin{quote}
\textit{To what extent is the predictive relation between asymmetric volatility and future stock returns driven by systematic versus idiosyncratic downside risk, and how does this relation vary with market conditions and trading frictions?}
\end{quote}
 
\noindent Two sub-questions structure the empirical analysis. The first asks whether RES contains predictive information for future stock returns beyond RSJ, which tests whether the RSJ effect is linked to left-tail severity or reflects a distinct form of return asymmetry. The second asks whether the predictive power of RSJ and RES is driven by systematic or idiosyncratic components, which directly tests the risk-premium versus mispricing interpretation.
 
This study connects the realized-volatility setting of \citet{bollerslev2020} with the high-frequency tail-risk framework of \citet{Dimitriadis03072022} in a unified cross-sectional design. By decomposing both measures and examining their predictive content across crisis states and arbitrage-friction conditions, it provides a structured empirical framework for assessing the relative importance of systematic risk compensation and firm-specific mispricing in explaining the cross-sectional relation between asymmetric volatility and future stock returns.
 
\section{Literature Review}

\subsection{Realized Asymmetric Volatility and Return Asymmetry}

The informational content of intraday return variation extends well beyond the level of volatility. \citet{barndorff2010} show that realized variance can be decomposed into positive and negative semivariance, and that negative semivariance has stronger predictive power for future volatility than its positive counterpart. \citet{bollerslev2022} extends this insight from forecasting to asset pricing, demonstrating that realized semi(co)variation measures carry distinct information for the pricing of equity returns. The central implication is that treating all volatility as equal obscures economically relevant variation in the sign structure of intraday returns.

This asymmetry has direct cross-sectional implications. \citet{amaya2015} show that realized skewness, constructed from high-frequency data, predicts the cross-section of equity returns: stocks with more negative realized skewness earn higher subsequent returns. \citet{bollerslev2020} build on this by introducing the Relative Signed Jump measure, defined as the normalized difference between positive and negative semivariance. RSJ is closely related to realized skewness but focuses specifically on the asymmetry in the volatility composition rather than the third moment of the return distribution. Importantly, \citet{bollerslev2020} show that RSJ subsumes the predictive power of realized skewness, volatility, and kurtosis in the cross-section, establishing it as the dominant higher-moment predictor of future stock returns. However, their analysis does not determine whether this predictability originates from systematic or firm-specific variation. The present study addresses this gap by decomposing RSJ into systematic and idiosyncratic components.

\subsection{Downside Tail Risk and Realized Expected Shortfall}

A related question is whether extreme negative outcomes, beyond the asymmetry captured by RSJ, are independently priced in the cross-section. \citet{artzner1999} establish Expected Shortfall as a coherent tail risk measure that captures the average severity of losses beyond a given quantile, rather than only a threshold level.

Empirical work confirms that downside risk is priced, with a progression in the literature toward higher-frequency measurement. Using monthly data, \citet{ang2006downside} show that downside beta, which measures a stock's comovement with the market conditional on market declines, is priced even after controlling for unconditional market beta. Using daily data, \citet{kelly2014} construct a time-varying tail risk factor from the cross-section of individual stock crashes and show that stocks with high tail risk loadings earn approximately 5\% higher annualized alpha. Also at the daily frequency, \citet{atilgan2020} document that stocks with recent extreme negative returns earn higher subsequent returns, attributing this to investor underreaction combined with costly arbitrage. Despite this progress, daily measures remain limited in their ability to capture within-day tail dynamics. \citet{Dimitriadis03072022} take the next step by constructing realized expected shortfall from intraday returns under a self-similar scaling law, producing more accurate tail risk estimates during periods of market stress. However, they apply RES to forecasting rather than to cross-sectional asset pricing.

The present study extends their framework to a cross-sectional setting, testing whether high-frequency tail severity predicts future stock returns. By analyzing RES alongside RSJ, it examines whether tail severity and volatility asymmetry capture overlapping or distinct dimensions of downside risk.

\subsection{Systematic versus Idiosyncratic Risk Decomposition}

Standard asset pricing theory implies that only systematic risk should be priced \citep{sharpe1964}. Yet empirical evidence on volatility-related measures presents a more complex picture. \citet{ang2006} document a puzzling negative relation between idiosyncratic volatility and subsequent returns, a finding that \citet{bali2008} show is sensitive to the treatment of small and illiquid stocks. For downside risk specifically, \citet{ang2006downside} demonstrate that it is the systematic component that commands a premium: downside beta is priced even after controlling for unconditional beta. These findings suggest that the cross-sectional effects of risk measures can differ substantially depending on whether the underlying variation is systematic or firm-specific.

The decomposition of realized asymmetric volatility into these components has received little attention in the literature. This is consequential, because the two components carry different economic implications. If a stock exhibits negative RSJ because its returns covary with market-wide downside movements, this systematic exposure should command a risk premium. If instead the negative RSJ reflects firm-specific shocks, standard theory predicts no compensation, unless limits to arbitrage prevent the mispricing from being corrected. The present study implements this decomposition at the intraday level using a high-frequency market model, applied to both RSJ and RES. For RSJ, two alternative methods are compared: one that separates intraday returns before constructing RSJ, and one that decomposes the final RSJ series using a market-wide RSJ factor.

\subsection{Limits to Arbitrage and Market Frictions}

If return predictability is driven by idiosyncratic components, the question is why it persists. The limits-to-arbitrage literature argues that arbitrage is risky and costly in practice, allowing mispricings to survive, particularly among stocks where trading frictions are most severe \citep{shleifer1997}. Empirical evidence supports this across multiple dimensions. \citet{amihud2002} shows that illiquidity is positively related to expected returns in the cross-section. \citet{bali2008} show that the cross-sectional effects of idiosyncratic volatility concentrate among stocks with the highest frictions, and the negative IVOL-return relation documented by \citet{ang2006} has been linked to the cost that idiosyncratic risk imposes on arbitrageurs. Together, these findings suggest that the persistence of volatility-related anomalies depends critically on the friction environment.

This motivates a key test in the present study. If the predictive power of the idiosyncratic components of RSJ and RES concentrates among stocks with high Amihud illiquidity, this supports a mispricing interpretation in which firm-specific downside risk persists because arbitrage frictions prevent its correction. If the systematic components remain predictive regardless of frictions, this favors a risk-premium explanation. This interaction-based test design follows the approach of \citet{atilgan2020} and extends it to the decomposed measures introduced in this study.

Taken together, the literature reviewed above converges on a common gap. Existing work establishes that asymmetric volatility predicts returns and that downside tail risk is priced, but these findings rely on aggregate measures that do not distinguish systematic from firm-specific sources. Conversely, the literature that does make this distinction has not been applied to realized semivariance-based or high-frequency tail risk measures. The present study bridges these approaches in a unified empirical framework.

\section{Data}

The empirical analysis uses five-minute intraday return data for U.S.\ equities provided
by the course instructor. The data are stored in daily files containing a vector of 78
five-minute discrete returns (\texttt{returns\_5m}) per stock per trading day, covering
the regular session from 9:30\textsc{am} to 4:00\textsc{pm} EST, alongside total daily
returns including dividends (\texttt{ret\_crsp}) and the number of intraday transactions
(\texttt{n\_obs}). Each stock is identified by a PERMNO, which is used to merge the
intraday data with CRSP to obtain weekly returns, market capitalizations, and firm-level
controls. The sample is restricted to ordinary common shares with CRSP share codes 10 or
11 listed on NYSE, AMEX, or NASDAQ, and to stocks with prices of at least \$5 to exclude
the most illiquid and micro-cap securities, following \citet{bollerslev2020}. All
five-minute discrete returns are converted to log returns prior to constructing the
realized measures.

The sample spans January 1993 through December 2024, extending the 1993--2013 period of
\citet{bollerslev2020} by approximately eleven years and covering major stress episodes
including the Global Financial Crisis and the COVID-19 pandemic. Robustness results
restricted to the original 1993--2013 period are reported in
Table~\ref{tab:sumstats_bollerslev} of Appendix~\ref{app:tables} to facilitate direct
comparison with the original study.

To reduce the influence of illiquid trading days on the realized measures, we apply a
liquidity screen similar in spirit to \citet{amaya2015}. A stock-day is classified as
valid only if at least 80 intraday transactions are recorded, i.e.,
$\texttt{n\_obs}_{i,d} \geq 80$, ensuring that the five-minute grid reflects
sufficiently active price formation. When a valid trading day contains isolated missing
five-minute intervals, those returns are set to zero, corresponding to the assumption that
no price change occurred during that interval. Weekly measures are constructed only for
stock-weeks in which at least three valid trading days pass this screen, so that the
realized measures rest on an adequate intraday sample. The resulting weekly cross-sectional
coverage over the full sample period is illustrated in Figure~\ref{fig:coverage} of
Appendix~\ref{app:figures}.

All intraday measures are constructed from information available up to the end of week $w$
and used to predict returns in week $w+1$. Weekly returns are measured from Tuesday close
to Tuesday close, following \citet{amaya2015}, which reduces microstructure distortions
associated with weekend return patterns. The intraday market proxy required for the
systematic decompositions of RSJ and RES is constructed as a value-weighted
cross-sectional average of five-minute returns across all available stocks in each
interval. As a robustness check, we replace this with five-minute returns on the SPDR
S\&P 500 ETF Trust (SPY), provided in a separate dataset with an identical structure. The
rolling coefficients used in the intraday market model are estimated over the four
preceding weeks of valid five-minute log returns and held fixed within week $w$. The
second RSJ decomposition, which operates at the weekly frequency, uses a rolling window
of 52 weeks.

Daily and monthly factor returns for the Fama--French--Carhart four-factor model are
obtained from Kenneth R.\ French's data library. Book values for the book-to-market ratio
are sourced from Compustat Fundamentals Annual, and stock returns are adjusted for
delisting using the CRSP delisting return to mitigate survivorship bias. The firm-level
control variables follow \citet{bollerslev2020} and include firm size (\textsc{me}),
book-to-market (\textsc{bm}), medium-term momentum (\textsc{mom}) \citep{jegadeesh1993},
short-term reversal (\textsc{rev}) \citep{lehmann1990, jegadeesh1990}, idiosyncratic
volatility (\textsc{ivol}) \citep{ang2006}, and the \citet{amihud2002} illiquidity
measure (\textsc{illiq}). All control variables are winsorized at the 1st and 99th
percentiles of their weekly cross-sectional distributions to limit the influence of
outliers. Pairwise correlations among all predictor variables are displayed in
Figure~\ref{fig:corr_heatmap} of Appendix~\ref{app:figures}.

The final panel, requiring non-missing values for the forward weekly return, RSJ, and all
baseline controls, contains approximately 3.8 million stock-week observations over the
full sample and approximately 1.9 million over the 1993--2013 subsample. Summary
statistics for all variables are provided in Table~\ref{tab:sumstats_full} of
Appendix~\ref{app:tables}; the corresponding statistics for the Bollerslev-comparable
subsample are reported in Table~\ref{tab:sumstats_bollerslev}.




\section{Methodology}
\subsection{Baseline RSJ Construction}

Following \citet{bollerslev2020}, let \(r_{i,t,k}\) denote the \(k\)-th
intraday return of stock \(i\) on day \(t\). Daily realized variance is

\begin{equation}
RV_{i,t} = \sum_{k=1}^{n} r_{i,t,k}^{2}.
\label{eq:rv}
\end{equation}

Positive and negative semivariance are

\begin{equation}
RV^{+}_{i,t} = \sum_{k=1}^{n} r_{i,t,k}^{2}\mathbf{1}(r_{i,t,k}>0),
\qquad
RV^{-}_{i,t} = \sum_{k=1}^{n} r_{i,t,k}^{2}\mathbf{1}(r_{i,t,k}<0).
\label{eq:semivar}
\end{equation}

The Relative Signed Jump measure is then

\begin{equation}
RSJ_{i,t} = \frac{RV^{+}_{i,t} - RV^{-}_{i,t}}{RV_{i,t}}.
\label{eq:rsj}
\end{equation}

Daily RSJ is aggregated to the weekly level by averaging across trading days
within week \(w\).

\subsection{Realized Expected Shortfall Construction}

Following \citet{Dimitriadis03072022}, we construct realized expected
shortfall directly from high-frequency returns under the
self-similar scaling framework. In the original paper, this construction
is implemented at the daily horizon. Because our empirical analysis is
conducted at the weekly frequency, we apply the same logic directly at
the weekly horizon by pooling all 5-minute returns within a week and
constructing weekly RES from that intraday sample.

Our raw data contain 5-minute simple returns. To align the construction
with \citet{Dimitriadis03072022}, who work with log-price increments, we
first transform each 5-minute simple return into a 5-minute log return.
Let \(R_{i,d,k}^{(5m)}\) denote the \(k\)-th 5-minute simple return of
stock \(i\) on day \(d\), for \(k=1,\dots,78\). We define the
corresponding 5-minute log return as
\begin{equation}
r_{i,d,k}^{(5m)}=\log\!\left(1+R_{i,d,k}^{(5m)}\right).
\label{eq:logret}
\end{equation}

To reduce the influence of illiquid stock-days, we apply a liquidity
screen analogous in spirit to \citet{amaya2015}. Specifically, we
classify day \(d\) for stock \(i\) as valid if the stock records at
least 80 intraday transactions on that day, that is,
\begin{equation}
n\_obs_{i,d}\ge 80.
\label{eq:screen}
\end{equation}
Accordingly, the daily decomposition files retain only stock-days that
satisfy this intraday liquidity screen. Let
\begin{equation}
D^{80}_{i,w}=\{d\in D_{i,w}: n\_obs_{i,d}\ge 80\}
\label{eq:validdays}
\end{equation}
denote the set of such valid trading days for stock \(i\) in week
\(w\). Weekly RES is then constructed by pooling the 5-minute log
returns from these retained valid days only. As RES is constructed
directly at the weekly horizon from pooled 5-minute returns, we require
at least three valid days in a week to compute weekly RES. Thus, weekly
RES is constructed only if
\begin{equation}
|D^{80}_{i,w}|\ge 3.
\label{eq:mindays}
\end{equation}

For each stock \(i\) and week \(w\), define the pooled vector of valid
5-minute log returns by
\begin{equation}
r_{i,w}^{c}
=
\left\{
r_{i,d,k}^{(5m)}: d\in D^{80}_{i,w},\; k=1,\dots,78
\right\}.
\label{eq:pooled}
\end{equation}
By construction, the pooled weekly return vector therefore contains only
5-minute log returns from stock-days that pass the daily intraday
liquidity screen. The total number of intraday observations entering the
weekly measure is therefore
\begin{equation}
c_{i,w}=78\,|D^{80}_{i,w}|.
\label{eq:count}
\end{equation}

For a given tail probability \(p\in(0,1)\), let
\(\widehat{Q}_{p}(r_{i,w}^{c})\) denote the empirical \(p\)-quantile of
the weekly intraday return sample, and let
\(\widehat{ES}_{p}(r_{i,w}^{c})\) denote the corresponding empirical
expected shortfall, that is, the sample average of 5-minute log returns
falling below the empirical lower-tail quantile:
\begin{equation}
\widehat{ES}_{p}(r_{i,w}^{c})
=
\frac{1}{N_{i,w}^{p}}
\sum_{r\in r_{i,w}^{c}}
r\,
\mathbf{1}\!\left(r\le \widehat{Q}_{p}(r_{i,w}^{c})\right),
\label{eq:es}
\end{equation}
where
\begin{equation}
N_{i,w}^{p}
=
\sum_{r\in r_{i,w}^{c}}
\mathbf{1}\!\left(r\le \widehat{Q}_{p}(r_{i,w}^{c})\right)
\label{eq:tailcount}
\end{equation}
is the number of weekly intraday observations in the lower tail.

Under the self-similar scaling framework of
\citet{Dimitriadis03072022}, the lower-frequency tail measure is
obtained by scaling the intraday expected shortfall by \(c_{i,w}^{H}\),
where \(H\) denotes the scaling exponent. Weekly realized expected
shortfall is therefore defined as
\begin{equation}
RES_{i,w}^{p}
=
c_{i,w}^{H}\,\widehat{ES}_{p}(r_{i,w}^{c}).
\label{eq:res}
\end{equation}

In the baseline specification, we set \(p=0.025\) and \(H=0.5\),
following \citet{Dimitriadis03072022}. Because RES is based on the
lower-tail mean of returns, it is typically negative in periods of
elevated downside risk. In the empirical analysis, we use the sign-adjusted measure
\begin{equation}
RES_{i,w}^{+,p}=-RES_{i,w}^{p},
\label{eq:res_adj}
\end{equation}
so that larger values correspond to greater downside tail risk.

\subsection{Systematic and Idiosyncratic Realized Expected Shortfall}

To distinguish market-wide from firm-specific tail risk, we extend the
realized expected shortfall framework by decomposing 5-minute stock
returns into systematic and idiosyncratic components using an intraday
market model with an intercept. We use 5-minute returns on the SPDR
S\&P 500 ETF Trust (SPY) as a proxy for the aggregate U.S. equity
market.

Let \(R_{SPY,d,k}^{(5m)}\) denote the 5-minute simple return on SPY in
interval \(k\) of day \(d\), and define the corresponding 5-minute log
return as
\begin{equation}
r_{SPY,d,k}^{(5m)}=\log\!\left(1+R_{SPY,d,k}^{(5m)}\right).
\label{eq:spy_logret}
\end{equation}

For stock \(i\), week \(w\), valid trading day \(d\in D^{80}_{i,w}\),
and intraday interval \(k=1,\dots,78\), we estimate the intraday market
model
\begin{equation}
r_{i,d,k}^{(5m)}
=
\hat{\alpha}_{i,w-1}
+
\hat{\beta}_{i,w-1}r_{SPY,d,k}^{(5m)}
+
\hat{\varepsilon}_{i,d,k},
\label{eq:market_model}
\end{equation}
where \(\hat{\alpha}_{i,w-1}\) and \(\hat{\beta}_{i,w-1}\) are estimated
from a rolling window of 5-minute log returns ending in week \(w-1\),
and are held fixed within week \(w\).

The fitted systematic component and residual idiosyncratic component are
then given by
\begin{equation}
r_{i,d,k}^{sys}
=
\hat{\alpha}_{i,w-1}
+
\hat{\beta}_{i,w-1}r_{SPY,d,k}^{(5m)},
\qquad
r_{i,d,k}^{idio}
=
\hat{\varepsilon}_{i,d,k}.
\label{eq:decomp_ret}
\end{equation}

The daily decomposition files retain only stock-days satisfying the
intraday liquidity screen \(n\_obs_{i,d}\ge 80\). Weekly systematic and
idiosyncratic RES are therefore constructed by pooling the corresponding
5-minute components only over retained valid days. As in the baseline
RES construction, we require at least three valid days in week \(w\),
that is, \(|D^{80}_{i,w}|\ge 3\).

Define
\begin{equation}
r_{i,w}^{sys,c}
=
\left\{
r_{i,d,k}^{sys}: d\in D^{80}_{i,w},\; k=1,\dots,78
\right\},
\qquad
r_{i,w}^{idio,c}
=
\left\{
r_{i,d,k}^{idio}: d\in D^{80}_{i,w},\; k=1,\dots,78
\right\}.
\label{eq:pooled_decomp}
\end{equation}

For each component, we compute the empirical lower-tail expected
shortfall analogously to the baseline RES measure:
\begin{equation}
\widehat{ES}_{p}(r_{i,w}^{sys,c})
=
\frac{1}{N_{i,w}^{sys,p}}
\sum_{r\in r_{i,w}^{sys,c}}
r\,
\mathbf{1}\!\left(r\le \widehat{Q}_{p}(r_{i,w}^{sys,c})\right),
\label{eq:es_sys}
\end{equation}
where
\begin{equation}
N_{i,w}^{sys,p}
=
\sum_{r\in r_{i,w}^{sys,c}}
\mathbf{1}\!\left(r\le \widehat{Q}_{p}(r_{i,w}^{sys,c})\right),
\label{eq:tailcount_sys}
\end{equation}
and
\begin{equation}
\widehat{ES}_{p}(r_{i,w}^{idio,c})
=
\frac{1}{N_{i,w}^{idio,p}}
\sum_{r\in r_{i,w}^{idio,c}}
r\,
\mathbf{1}\!\left(r\le \widehat{Q}_{p}(r_{i,w}^{idio,c})\right),
\label{eq:es_idio}
\end{equation}
where
\begin{equation}
N_{i,w}^{idio,p}
=
\sum_{r\in r_{i,w}^{idio,c}}
\mathbf{1}\!\left(r\le \widehat{Q}_{p}(r_{i,w}^{idio,c})\right).
\label{eq:tailcount_idio}
\end{equation}

Weekly systematic and idiosyncratic realized expected shortfall are then
defined as
\begin{equation}
RES_{i,w}^{sys,p}
=
c_{i,w}^{H}\,\widehat{ES}_{p}(r_{i,w}^{sys,c}),
\qquad
RES_{i,w}^{idio,p}
=
c_{i,w}^{H}\,\widehat{ES}_{p}(r_{i,w}^{idio,c}).
\label{eq:res_decomp}
\end{equation}

For ease of interpretation, we use the sign-adjusted versions
\begin{equation}
RES_{i,w}^{sys+,p}=-RES_{i,w}^{sys,p},
\qquad
RES_{i,w}^{idio+,p}=-RES_{i,w}^{idio,p},
\label{eq:res_decomp_adj}
\end{equation}
so that larger values correspond to greater downside systematic and
idiosyncratic tail risk, respectively.

\subsection{Two Alternative Decompositions of RSJ}

We consider two alternative decompositions of RSJ. The baseline method,
Method 1, decomposes intraday returns before RSJ construction, whereas
Method 2 decomposes the final RSJ series directly. Comparing both
methods allows us to assess whether the results depend on the
decomposition approach.

While both methods are implemented and compared, Method 1 is preferred
as the primary decomposition for three reasons. First, it operates at
the level of the underlying intraday returns before RSJ is constructed,
so that the systematic component directly reflects the asymmetry in the
market-driven part of intraday price variation. This makes the economic
interpretation more transparent. Second, Method 1 is methodologically
consistent with the decomposition of realized expected shortfall
presented below, which is also implemented at the intraday return
level. This consistency allows for a cleaner comparison between the
systematic and idiosyncratic components of RSJ and RES. Third, because
RSJ is a nonlinear ratio, decomposing it only after construction, as in
Method 2, implies that the decomposition is applied to an already
aggregated measure, thereby introducing an additional layer of
aggregation that is absent in Method 1.

For these reasons, the results based on Method 1 are treated as the
primary findings, while Method 2 serves as a robustness check.
Discrepancies between the two methods are interpreted as evidence that
the decomposition is method-sensitive, in which case the conclusions
based on Method 1 are given precedence.

\subsubsection{Method 1: Decomposition Before RSJ Construction}

In the first method, intraday returns are decomposed into systematic and
idiosyncratic components before realized semivariance and RSJ are
constructed. Using the same high-frequency market model as in the RES
decomposition, for stock \(i\), trading day \(d \in D_w\), and intraday
interval \(k\), we write

\begin{equation}
r_{i,d,k} = \hat{\alpha}_{i,w-1} + \hat{\beta}_{i,w-1} r_{M,d,k} + \hat{\varepsilon}_{i,d,k},
\label{eq:m1_model}
\end{equation}

where \(r_{M,d,k}\) denotes the intraday market return. The coefficients
\(\hat{\alpha}_{i,w-1}\) and \(\hat{\beta}_{i,w-1}\) are estimated using a
rolling window of intraday returns ending in week \(w-1\), and are held
fixed within week \(w\). In the baseline specification, we use the
previous four weeks of 5-minute returns. This yields

\begin{equation}
r^{sys}_{i,d,k} = \hat{\alpha}_{i,w-1} + \hat{\beta}_{i,w-1} r_{M,d,k},
\qquad
r^{idio}_{i,d,k} = \hat{\varepsilon}_{i,d,k}.
\label{eq:m1_decomp}
\end{equation}

Using these two intraday return series, realized semivariance is then
computed separately for the systematic and idiosyncratic components:

\begin{equation}
RV_{i,d}^{+,\mathrm{sys}} =
\sum_k \left(r_{i,d,k}^{\mathrm{sys}}\right)^2
\mathbf{1}(r_{i,d,k}^{\mathrm{sys}} > 0),
\qquad
RV_{i,d}^{-,\mathrm{sys}} =
\sum_k \left(r_{i,d,k}^{\mathrm{sys}}\right)^2
\mathbf{1}(r_{i,d,k}^{\mathrm{sys}} < 0),
\label{eq:semivar_sys}
\end{equation}

\begin{equation}
RV_{i,d}^{+,\mathrm{idio}} =
\sum_k \left(r_{i,d,k}^{\mathrm{idio}}\right)^2
\mathbf{1}(r_{i,d,k}^{\mathrm{idio}} > 0),
\qquad
RV_{i,d}^{-,\mathrm{idio}} =
\sum_k \left(r_{i,d,k}^{\mathrm{idio}}\right)^2
\mathbf{1}(r_{i,d,k}^{\mathrm{idio}} < 0).
\label{eq:semivar_idio}
\end{equation}

Daily systematic and idiosyncratic RSJ are then defined as

\begin{equation}
RSJ_{i,d}^{sys}
=
\frac{RV_{i,d}^{+,\mathrm{sys}} - RV_{i,d}^{-,\mathrm{sys}}}
{RV_{i,d}^{+,\mathrm{sys}} + RV_{i,d}^{-,\mathrm{sys}}},
\qquad
RSJ_{i,d}^{idio}
=
\frac{RV_{i,d}^{+,\mathrm{idio}} - RV_{i,d}^{-,\mathrm{idio}}}
{RV_{i,d}^{+,\mathrm{idio}} + RV_{i,d}^{-,\mathrm{idio}}}.
\label{eq:rsj_daily}
\end{equation}

To match the baseline RSJ construction, these daily measures are
aggregated to the weekly frequency by averaging across trading days in
week \(w\):

\begin{equation}
RSJ_{i,w}^{sys}
=
\frac{1}{N_w}\sum_{d \in D_w} RSJ_{i,d}^{sys},
\qquad
RSJ_{i,w}^{idio}
=
\frac{1}{N_w}\sum_{d \in D_w} RSJ_{i,d}^{idio}.
\label{eq:rsj_weekly_m1}
\end{equation}

Because RSJ is a nonlinear ratio, the resulting measures should be
interpreted as model-based decompositions rather than exact algebraic
components of the original RSJ. Nevertheless, they provide an
economically meaningful distinction between asymmetric systematic
volatility and asymmetric firm-specific volatility.

\subsubsection{Method 2: Decomposition of RSJ After Construction}

In the second method, firm-level RSJ is first constructed in the standard
way and then decomposed using an aggregate market-wide RSJ factor:

\begin{equation}
RSJ_{i,w} = \hat{\alpha}_{i,w-1} + \hat{\beta}_{i,w-1}^{RSJ} RSJ^{M}_w + \hat{\varepsilon}_{i,w},
\label{eq:m2_reg}
\end{equation}

where \(RSJ^{M}_w\) denotes the market-wide RSJ measure, constructed as a
value-weighted cross-sectional average of firm-level RSJ. The coefficients
\(\hat{\alpha}_{i,w-1}\) and \(\hat{\beta}_{i,w-1}^{RSJ}\) are estimated
using a rolling window of weekly observations ending in week \(w-1\). In
the baseline specification, we use the previous 52 weeks of data so that
only information available prior to week \(w\) enters the decomposition.

The fitted value

\begin{equation}
RSJ^{sys}_{i,w}
=
\hat{\alpha}_{i,w-1}
+
\hat{\beta}_{i,w-1}^{RSJ} RSJ^{M}_w
\label{eq:rsj_sys_m2}
\end{equation}

captures the systematic component, while the residual

\begin{equation}
RSJ^{idio}_{i,w}
=
\hat{\varepsilon}_{i,w}
\label{eq:rsj_idio_m2}
\end{equation}

captures the idiosyncratic component.

Under this definition, the systematic component corresponds to the fitted
part of the market-model regression, so that the decomposition is exact
in a statistical sense, although the intercept is not itself market-driven.
This second method has the advantage of directly decomposing the final
RSJ measure used in the prediction tests. Comparing both methods allows
us to assess whether the source of predictability depends on the
decomposition approach. Because Method 1 decomposes intraday returns
prior to RSJ construction, it provides a more direct economic
distinction between systematic and firm-specific asymmetric volatility.
Method 2 instead decomposes the final RSJ measure itself and therefore
serves as an alternative specification. Divergence between the two
methods would indicate that the decomposition of RSJ is sensitive to
implementation, which we treat as an empirical finding rather than a
limitation.

\subsection{Portfolio Sorting Procedure}

To evaluate the predictive power of the proposed risk measures, firms are
sorted into portfolios based on their estimated RSJ and RES measures. At
the end of each week \(w\), all stocks in the sample are ranked according
to the cross-sectional distribution of the relevant variable and assigned
to quintile portfolios.

Portfolio \(Q1\) contains stocks with the lowest values of the sorting
variable, while portfolio \(Q5\) contains stocks with the highest values.
The sorting procedure is repeated separately for each of the following
variables: \(RSJ_{i,w}\), \(RSJ^{sys}_{i,w}\), \(RSJ^{idio}_{i,w}\),
\(RES^{\alpha}_{i,w}\), \(RES^{sys,\alpha}_{i,w}\), and
\(RES^{idio,\alpha}_{i,w}\).

After portfolio assignment, returns are measured over the following week.
For each portfolio, we compute both equal-weighted and value-weighted
returns. The equal-weighted return on portfolio \(Qj\) in week \(w+1\) is

\begin{equation}
R^{EW}_{Qj,w+1}
=
\frac{1}{N_{Qj,w}}
\sum_{i \in Qj} R_{i,w+1},
\label{eq:ew_return}
\end{equation}

where \(N_{Qj,w}\) denotes the number of firms in portfolio \(Qj\) during
week \(w\). The value-weighted return is defined as

\begin{equation}
R^{VW}_{Qj,w+1}
=
\sum_{i \in Qj} \omega_{i,w} R_{i,w+1},
\label{eq:vw_return}
\end{equation}

where \(\omega_{i,w}\) denotes the weight of stock \(i\) in portfolio
\(Qj\) based on its market capitalization at the end of week \(w\), with
\(\sum_{i \in Qj} \omega_{i,w} = 1\).

To assess the economic magnitude of predictability, we construct spread
portfolios that are aligned with the theoretical prediction for each risk
measure. For RES-based sorts, where higher downside tail risk is expected
to be associated with higher future returns, we compute high-minus-low
spreads:

\begin{equation}
R^{EW,H-L}_{w+1} = R^{EW}_{Q5,w+1} - R^{EW}_{Q1,w+1},
\qquad
R^{VW,H-L}_{w+1} = R^{VW}_{Q5,w+1} - R^{VW}_{Q1,w+1}.
\label{eq:spread_hml}
\end{equation}

For RSJ-based sorts, where lower (more negative) RSJ is expected to be
associated with higher future returns, we compute low-minus-high spreads:

\begin{equation}
R^{EW,L-H}_{w+1} = R^{EW}_{Q1,w+1} - R^{EW}_{Q5,w+1},
\qquad
R^{VW,L-H}_{w+1} = R^{VW}_{Q1,w+1} - R^{VW}_{Q5,w+1}.
\label{eq:spread_lmh}
\end{equation}

A positive and statistically significant spread therefore indicates that
the portfolio results are consistent with the predicted direction of the
corresponding risk measure.


\subsection{Cross-Sectional Predictive Regressions}

The empirical analysis is conducted at the stock-week level. For each
stock~$i$ and week~$w$, we construct weekly measures of asymmetric
volatility and tail risk from intraday data observed up to the end of
week~$w$, and relate these measures to stock returns in week~$w+1$.
We follow the Fama--MacBeth methodology: for each week~$w$ we estimate
a cross-sectional regression, average the slope coefficients over time,
and assess statistical significance using Newey--West heteroskedasticity
and autocorrelation consistent standard errors applied to the resulting
time series of weekly estimates.

Throughout, $X_{i,w}$ denotes the vector of firm-level controls,
comprising log market equity, book-to-market ratio, 12-2 month
momentum, one-week reversal, idiosyncratic volatility, and the
Amihud~(2002) illiquidity measure. For RSJ-based measures, the
predicted sign of the slope coefficient is negative, reflecting the
cross-sectional relation between more negative RSJ and higher future
returns. For RES-based measures, the predicted sign is positive,
reflecting that greater downside tail risk is associated with higher
subsequent returns.

\paragraph{Aggregate specifications.}
We begin with three specifications that use the aggregate, non-decomposed
measures. Specification~B1 regresses next-week returns on RSJ alone
with controls:
\begin{equation}
R_{i,w+1} = \alpha_w + \beta_{1,w}\,RSJ_{i,w} + \gamma_w' X_{i,w}
+ \varepsilon_{i,w+1}.
\label{eq:B1}
\end{equation}
Specification~B2 replaces RSJ with RES:
\begin{equation}
R_{i,w+1} = \alpha_w + \beta_{1,w}\,RES_{i,w} + \gamma_w' X_{i,w}
+ \varepsilon_{i,w+1}.
\label{eq:B2}
\end{equation}
Specification~B3 includes both measures jointly:
\begin{equation}
R_{i,w+1} = \alpha_w + \beta_{1,w}\,RSJ_{i,w} + \beta_{2,w}\,RES_{i,w}
+ \gamma_w' X_{i,w} + \varepsilon_{i,w+1}.
\label{eq:B3}
\end{equation}
These three specifications establish the baseline predictive content of
each measure and test whether RES retains incremental information beyond
RSJ when both are included simultaneously.

\paragraph{RSJ decomposition specifications.}
To examine whether the predictive content of RSJ originates from
systematic or idiosyncratic asymmetric volatility, we estimate
decomposition specifications separately for each decomposition method.
Specification~B4, which serves as the primary decomposition, uses
Method~1 and includes both RSJ components alongside controls:
\begin{equation}
R_{i,w+1} = \alpha_w + \beta_{1,w}\,RSJ^{sys,M1}_{i,w}
+ \beta_{2,w}\,RSJ^{idio,M1}_{i,w}
+ \gamma_w' X_{i,w} + \varepsilon_{i,w+1}.
\label{eq:B4}
\end{equation}
Specification~B5 repeats this using Method~2 as a robustness check:
\begin{equation}
R_{i,w+1} = \alpha_w + \beta_{1,w}\,RSJ^{sys,M2}_{i,w}
+ \beta_{2,w}\,RSJ^{idio,M2}_{i,w}
+ \gamma_w' X_{i,w} + \varepsilon_{i,w+1}.
\label{eq:B5}
\end{equation}

\paragraph{RES decomposition specification.}
Specification~B6 isolates the systematic and idiosyncratic components of
realized expected shortfall:
\begin{equation}
R_{i,w+1} = \alpha_w + \beta_{1,w}\,RES^{sys}_{i,w}
+ \beta_{2,w}\,RES^{idio}_{i,w}
+ \gamma_w' X_{i,w} + \varepsilon_{i,w+1}.
\label{eq:B6}
\end{equation}

\paragraph{Full decomposition specifications.}
The main specifications of interest include all four decomposed
components simultaneously. Specification~B7 uses Method~1 for the RSJ
decomposition:
\begin{equation}
R_{i,w+1} = \alpha_w
+ \beta_{1,w}\,RSJ^{sys,M1}_{i,w}
+ \beta_{2,w}\,RSJ^{idio,M1}_{i,w}
+ \beta_{3,w}\,RES^{sys}_{i,w}
+ \beta_{4,w}\,RES^{idio}_{i,w}
+ \gamma_w' X_{i,w}
+ \varepsilon_{i,w+1},
\label{eq:main_fm}
\end{equation}
and specification~B8 repeats this using Method~2:
\begin{equation}
R_{i,w+1} = \alpha_w
+ \beta_{1,w}\,RSJ^{sys,M2}_{i,w}
+ \beta_{2,w}\,RSJ^{idio,M2}_{i,w}
+ \beta_{3,w}\,RES^{sys}_{i,w}
+ \beta_{4,w}\,RES^{idio}_{i,w}
+ \gamma_w' X_{i,w}
+ \varepsilon_{i,w+1}.
\label{eq:B8}
\end{equation}
Equation~\eqref{eq:main_fm} constitutes the primary empirical
specification of this paper. Comparing B7 and B8 allows us to assess
whether the results are sensitive to the choice of RSJ decomposition
method.

\paragraph{Idiosyncratic dominance specifications.}
Finally, two additional specifications test whether the idiosyncratic
channel is robust to the inclusion of its systematic counterpart.
Specification~B9 includes idiosyncratic RSJ (Method~1) and idiosyncratic
RES jointly:
\begin{equation}
R_{i,w+1} = \alpha_w
+ \beta_{1,w}\,RSJ^{idio,M1}_{i,w}
+ \beta_{2,w}\,RES^{idio}_{i,w}
+ \gamma_w' X_{i,w}
+ \varepsilon_{i,w+1}.
\label{eq:B9}
\end{equation}
Specification~B10 replaces idiosyncratic RSJ with its systematic
counterpart to examine whether systematic RSJ retains any predictive
power once idiosyncratic tail risk is controlled for:
\begin{equation}
R_{i,w+1} = \alpha_w
+ \beta_{1,w}\,RSJ^{sys,M1}_{i,w}
+ \beta_{2,w}\,RES^{idio}_{i,w}
+ \gamma_w' X_{i,w}
+ \varepsilon_{i,w+1}.
\label{eq:B10}
\end{equation}

The average risk-premium estimate for each slope is
\begin{equation}
\bar{\beta}_j = \frac{1}{W}\sum_{w=1}^{W} \hat{\beta}_{j,w},
\label{eq:fm_avg}
\end{equation}
where $W$ denotes the total number of weeks in the sample.

\subsection{Economic Interpretation}

The decomposition in equation~\eqref{eq:main_fm} allows us to distinguish between
systematic and idiosyncratic sources of downside-risk predictability.

A risk-based interpretation predicts that the systematic components,
$RSJ^{sys}_{i,w}$ and $RES^{sys}_{i,w}$, should be most strongly related to future returns.
A mispricing-based interpretation instead predicts that the idiosyncratic components,
$RSJ^{idio}_{i,w}$ and $RES^{idio}_{i,w}$, should contain the strongest predictive content.

In addition, including both RSJ and RES in the same regression allows us to examine whether
RES contains incremental predictive information beyond RSJ.

\subsection{Hypothesis Tests}

Let
\[
\bar{\bm{\beta}}
=
\begin{pmatrix}
\bar{\beta}_1\\
\bar{\beta}_2\\
\bar{\beta}_3\\
\bar{\beta}_4
\end{pmatrix}
\]
denote the vector of time-averaged slope coefficients from equation~\eqref{eq:main_fm}.
We evaluate four main hypotheses concerning the
relative importance of systematic and idiosyncratic downside-risk
components.

\paragraph{Test 1: Systematic versus idiosyncratic RSJ}
\begin{equation}
H_0^{(1)}:\ \bar{\beta}_1 = \bar{\beta}_2.
\label{eq:H1}
\end{equation}

\paragraph{Test 2: Systematic versus idiosyncratic RES}
\begin{equation}
H_0^{(2)}:\ \bar{\beta}_3 = \bar{\beta}_4.
\label{eq:H2}
\end{equation}

\paragraph{Test 3: Joint significance of idiosyncratic downside risk}
\begin{equation}
H_0^{(3)}:\ \bar{\beta}_2 = 0
\quad \text{and} \quad
\bar{\beta}_4 = 0.
\label{eq:H3}
\end{equation}

\paragraph{Test 4: Joint significance of systematic downside risk}
\begin{equation}
H_0^{(4)}:\ \bar{\beta}_1 = 0
\quad \text{and} \quad
\bar{\beta}_3 = 0.
\label{eq:H4}
\end{equation}

Tests (1) and (2) compare the systematic and idiosyncratic components of
RSJ and RES respectively. Tests (3) and (4) evaluate the joint
explanatory power of the idiosyncratic and systematic downside-risk
blocks. Comparing these two tests provides a direct assessment of
whether cross-sectional predictability is primarily driven by
firm-specific or market-wide downside risk.

\subsection{Wald Test Implementation}

The hypotheses above are evaluated using Wald tests based on the average
coefficient vector $\bar{\bm{\beta}}$ and its estimated covariance matrix
$\widehat{\mathrm{Var}}(\bar{\bm{\beta}})$.

For a set of linear restrictions
\[
H_0:\ R\bar{\bm{\beta}} = r,
\]
the Wald statistic is

\begin{equation}
W
=
(R\bar{\bm{\beta}} - r)'
\left[
R\,\widehat{\mathrm{Var}}(\bar{\bm{\beta}})\,R'
\right]^{-1}
(R\bar{\bm{\beta}} - r),
\label{eq:wald}
\end{equation}

which is asymptotically $\chi^2(q)$ under the null, where $q$ denotes
the number of restrictions.

The restriction matrices used in the four tests are

\begin{equation}
R_1 =
\begin{pmatrix}
1 & -1 & 0 & 0
\end{pmatrix},
\qquad
r_1 = 0,
\label{eq:R1}
\end{equation}

\begin{equation}
R_2 =
\begin{pmatrix}
0 & 0 & 1 & -1
\end{pmatrix},
\qquad
r_2 = 0,
\label{eq:R2}
\end{equation}

\begin{equation}
R_3 =
\begin{pmatrix}
0 & 1 & 0 & 0\\
0 & 0 & 0 & 1
\end{pmatrix},
\qquad
r_3 =
\begin{pmatrix}
0\\
0
\end{pmatrix},
\label{eq:R3}
\end{equation}

\begin{equation}
R_4 =
\begin{pmatrix}
1 & 0 & 0 & 0\\
0 & 0 & 1 & 0
\end{pmatrix},
\qquad
r_4 =
\begin{pmatrix}
0\\
0
\end{pmatrix}.
\label{eq:R4}
\end{equation}


\subsection{Crisis-State Analysis}

We next study whether the predictive effects of systematic downside risk
strengthen during crisis periods, defined in the baseline specification
using NBER recession dates. Specifically, we define a crisis-state
indicator that equals one during weeks that fall within officially dated
NBER recession periods, and zero otherwise.

Let \(C_w\) denote the recession-state dummy equal to one during NBER
recession weeks. This definition captures periods of broad
macroeconomic contraction and heightened financial stress, which are
natural environments in which systematic downside risk may become more
relevant for asset pricing.

State dependence is tested using a Fama--MacBeth interaction regression
that includes the systematic components of both RSJ and RES jointly:

\begin{equation}
\begin{aligned}
R_{i,w+1}
&=
\alpha_w
+
\beta_{1,w} RSJ^{sys}_{i,w}
+
\beta_{2,w} RES^{sys,\alpha}_{i,w}
+
\beta_{3,w} C_w \\
&\quad
+
\beta_{4,w} \bigl(RSJ^{sys}_{i,w}\times C_w\bigr)
+
\beta_{5,w} \bigl(RES^{sys,\alpha}_{i,w}\times C_w\bigr)
+
\gamma_w' X_{i,w}
+
\varepsilon_{i,w+1}.
\end{aligned}
\label{eq:crisis}
\end{equation}

where \(X_{i,w}\) denotes the standard vector of firm-level control
variables.

In this specification, \(\beta_{1,w}\) and \(\beta_{2,w}\) capture the
predictive effects of systematic RSJ and systematic RES during
non-recession weeks, while \(\beta_{4,w}\) and \(\beta_{5,w}\) measure
whether these predictive effects change during NBER recession periods.
A positive and significant interaction effect for systematic downside
risk would indicate that its pricing becomes stronger in crisis states.

If the predictive effects of \(RSJ^{sys}_{i,w}\) and
\(RES^{sys,\alpha}_{i,w}\) strengthen during recessions, this would
support a risk-based interpretation, since systematic downside exposures
should matter most when aggregate economic conditions deteriorate. By
contrast, weak or insignificant crisis-state amplification in these
systematic components would provide less support for a rational
risk-premium explanation.

\subsection{Amihud Illiquidity and Its Role in the Empirical Design}

We measure trading frictions using the Amihud (2002) illiquidity proxy.
For stock $i$ on day $d$, daily illiquidity is defined as
\begin{equation}
\textit{ILLIQ}^{daily}_{i,d}
=
\frac{|R^{adj}_{i,d}|}{\textit{DollarVol}_{i,d}},
\label{eq:illiq_daily}
\end{equation}
where $R^{adj}_{i,d}$ is the delisting-adjusted daily return and
\begin{equation}
\textit{DollarVol}_{i,d}=|P_{i,d}|\times \textit{Volume}_{i,d}.
\label{eq:dollarvol}
\end{equation}

To obtain a stable weekly friction measure, we first construct a rolling
5-trading-day average:
\begin{equation}
\widetilde{\textit{ILLIQ}}_{i,d}
=
\frac{1}{5}\sum_{j=0}^{4}\textit{ILLIQ}^{daily}_{i,d-j},
\label{eq:illiq_roll}
\end{equation}
requiring five valid daily observations. We then log-transform:
\begin{equation}
\textit{ILLIQ}_{i,d}=\log\!\left(\widetilde{\textit{ILLIQ}}_{i,d}\right),
\label{eq:illiq_log}
\end{equation}
and map this to the weekly panel by taking the end-of-week (Tuesday-week)
snapshot for each stock-week.

In the return-predictability regressions, $\textit{ILLIQ}_{i,w}$ enters as a
standard control alongside $ME$, $BM$, $MOM$, $REV$, and $IVOL$.
All controls, including $\textit{ILLIQ}_{i,w}$, are winsorized cross-sectionally
at the 1st and 99th percentiles each week before estimation.

\subsection{Illiquidity Mechanism Check}

As an additional descriptive mechanism test, we estimate weekly
cross-sectional regressions of idiosyncratic downside-risk characteristics
on illiquidity:
\begin{equation}
C_{i,w}
=
\alpha_w
+
\delta_w\,ILLIQ_{i,w}
+
\eta_w'X_{i,w}
+
\nu_{i,w},
\label{eq:mech}
\end{equation}
where \(C_{i,w}\) denotes one of the idiosyncratic downside-risk measures,
such as \(RSJ^{idio}_{i,w}\) or \(RES^{idio}_{i,w}\), and
\begin{equation}
X_{i,w}=(ME,BM,MOM,REV,IVOL).
\label{eq:controls_mech}
\end{equation}
Thus, illiquidity is the focal regressor and is not duplicated inside
\(X_{i,w}\). The reported coefficient is the time-series average of
\(\hat{\delta}_w\), with Newey--West standard errors (6 lags).

We report both univariate and controlled specifications. The univariate
regression captures the raw association between illiquidity and the
relevant downside-risk characteristic, while the controlled specification
shows whether this association remains after accounting for standard firm
characteristics. This mechanism exercise is estimated separately from the
main return-predictability Fama--MacBeth models and is interpreted as
descriptive evidence on exposure--friction comovement rather than as a
pricing interaction model.

\section{Results}
\label{sec:results}

The empirical evidence is organised in five parts, following the structure of the
methodology. Section~\ref{sec:portfolio_sorts} establishes the direction and economic
magnitude of return premia through quintile portfolio sorts. Section~\ref{sec:fm}
evaluates robustness to standard risk-factor controls via Fama--MacBeth cross-sectional
regressions and documents the complementarity between RSJ and RES. Section~\ref{sec:decomp}
decomposes both measures into systematic and idiosyncratic components, tests their
relative risk prices with formal Wald statistics, and shows that virtually all predictive
content is concentrated in firm-specific variation. Section~\ref{sec:crisis} examines
whether the systematic components are amplified during NBER recession periods, as a
rational risk-compensation view would require. Section~\ref{sec:frictions} tests whether
idiosyncratic predictability concentrates among illiquid, hard-to-arbitrage stocks.
Across all five analyses the evidence points consistently toward a limits-to-arbitrage
interpretation rather than rational risk compensation.

%%------------------------------------------------------------
\subsection{Portfolio Sorts}
\label{sec:portfolio_sorts}
%%------------------------------------------------------------

Each week, stocks are sorted into five equal-count quintile portfolios based on the
relevant signal. For RSJ-based signals a low-minus-high (L$-$H) spread is constructed;
for RES-based signals a high-minus-low (H$-$L) spread is constructed. Both equal-weighted
(EW) and value-weighted (VW) average weekly spread returns are reported. Table~\ref{tab:portfolio_spreads}
presents spread means and $t$-statistics.

\begin{table}[!htbp]
\centering
\caption{Quintile Portfolio Sorts: Spread Returns}
\label{tab:portfolio_spreads}
\small
\begin{threeparttable}
\begin{tabular}{l c c c c}
\toprule
Signal
  & \multicolumn{2}{c}{EW (bps/week)}
  & \multicolumn{2}{c}{VW (bps/week)} \\
\cmidrule(lr){2-3}\cmidrule(lr){4-5}
  & Spread & $t$-stat & Spread & $t$-stat \\
\midrule
\multicolumn{5}{l}{\textit{Panel A: RSJ-based signals, low-minus-high (L$-$H) spreads}} \\[3pt]
RSJ Total              & $23.49^{***}$ & $(5.85)$ & $17.55^{***}$ & $(3.51)$ \\
\quad RSJ Systematic (M1) & $-2.26$    & $(-0.51)$ & $3.01$       & $(0.54)$ \\
\quad RSJ Idiosyncratic (M1) & $17.41^{***}$ & $(4.14)$ & $13.14^{*}$ & $(2.35)$ \\
\quad RSJ Systematic (M2) & $-6.18$    & $(-1.30)$ & $0.37$       & $(0.07)$ \\
\quad RSJ Idiosyncratic (M2) & $19.57^{***}$ & $(5.53)$ & $14.88^{**}$ & $(3.10)$ \\[5pt]
\multicolumn{5}{l}{\textit{Panel B: RES-based signals, high-minus-low (H$-$L) spreads}} \\[3pt]
RES Total              & $1.90$ & $(0.20)$ & $1.55$ & $(0.17)$ \\
\quad RES Systematic   & $-0.18$ & $(-0.03)$ & $3.43$ & $(0.46)$ \\
\quad RES Idiosyncratic & $0.59$ & $(0.07)$ & $0.51$ & $(0.06)$ \\
\bottomrule
\end{tabular}
\begin{tablenotes}
\scriptsize
\item \textit{Notes:} At the end of each week, stocks are sorted into five equal-count
quintile portfolios based on the relevant signal. Panel~A reports low-minus-high (L$-$H)
spreads for RSJ-based signals (Q1$-$Q5); Panel~B reports high-minus-low (H$-$L) spreads
for RES-based signals (Q5$-$Q1). EW (VW) denotes equal-weighted (value-weighted) weekly
returns, averaged over time. $t$-statistics are computed with Newey--West standard errors
(six lags). The corresponding FF3 alphas (not tabulated here) follow the same directional
pattern as the raw spreads.
Sample: January 1993--December 2024. \\
$^{*}p<0.10$,\ $^{**}p<0.05$,\ $^{***}p<0.01$.
\end{tablenotes}
\end{threeparttable}
\end{table}

The RSJ total spread earns 23.49 basis points per week equal-weighted ($t=5.85$) and
17.55 basis points value-weighted ($t=3.51$), with FF3 alphas of 22.95 and 17.40 basis
points respectively, both significant at the 1\% level. These findings replicate and
extend \citet{bollerslev2020}, confirming that the negative cross-sectional relation
between RSJ and future stock returns holds robustly over the longer 1993--2024 sample.

The decomposition results reveal a striking pattern: all of the RSJ portfolio premium
resides in the idiosyncratic component. Under Method~1, the idiosyncratic spread earns
17.41 bps EW ($t=4.14$), and under Method~2 it reaches 19.57 bps ($t=5.53$). The
systematic component, by contrast, is economically negligible and statistically
insignificant under both methods (M1: $-2.26$ bps, $t=-0.51$; M2: $-6.18$ bps,
$t=-1.30$), and even switches sign. Under a risk-compensation view, the systematic
component should carry the predictive content. The opposite pattern --- large idiosyncratic
premium, near-zero systematic premium --- is directly consistent with a mispricing
mechanism driven by limits to arbitrage.

For RES, all raw long-short spreads are economically negligible and statistically
insignificant, across both weighting schemes and all three variants (total, systematic,
idiosyncratic). This does not imply that RES lacks return-predictive content; rather, it
reflects the fact that RES is strongly correlated with established predictors such as
firm size and idiosyncratic volatility. Section~\ref{sec:fm} shows that once these
covariates are controlled for in a regression framework, RES contains independent
predictive information.

Figure~\ref{fig:portfolio_spreads_fig} visualises the spread returns across all signals,
making the contrast between idiosyncratic and systematic RSJ immediately apparent. The
idiosyncratic RSJ bars are large and consistent across both decomposition methods; the
systematic RSJ bars sit near zero and flip sign depending on the method. The RES bars
cluster uniformly near zero, corroborating the statistical findings in Panel~B of
Table~\ref{tab:portfolio_spreads}.

\begin{figure}[!ht]
\centering
\includegraphics[width=0.88\textwidth]{results/figures/fig_portfolio_spreads_ew_bps}
\caption{Equal-weighted quintile portfolio spread returns (bps/week) for all RSJ and
RES signals. RSJ spreads are low-minus-high (L$-$H); RES spreads are high-minus-low
(H$-$L). Idiosyncratic RSJ generates large positive spreads under both M1 and M2,
comparable to the total RSJ spread, while systematic RSJ is near zero under both methods.
All RES long-short means are economically negligible.}
\label{fig:portfolio_spreads_fig}
\end{figure}

The persistence of the idiosyncratic RSJ premium over time is illustrated in
Figure~\ref{fig:cumulative}. The total RSJ and idiosyncratic RSJ (M2) long-short
strategies accumulate returns steadily throughout the 1993--2024 sample. Critically,
the post-2013 period --- which lies outside the original \citet{bollerslev2020} sample
--- shows no sign of premium attenuation, confirming that the finding is not confined to
the original estimation window. By contrast, the RES long-short portfolios remain
essentially flat throughout, underscoring that portfolio sorts cannot recover RES
predictability without controlling for firm characteristics.

\begin{figure}[!ht]
\centering
\includegraphics[width=0.88\textwidth]{results/figures/fig_cumulative_long_short_performance}
\caption{Cumulative equal-weighted long-short portfolio returns (in percent, weekly
rebalancing) over January 1993 to December 2024 for RSJ total (L$-$H), RSJ idiosyncratic
M2 (L$-$H), RES total (H$-$L), and RES idiosyncratic (H$-$L). The RSJ total and
idiosyncratic strategies accumulate returns persistently throughout the sample, including
the post-2013 period not covered by \citet{bollerslev2020}. RES long-short portfolios
remain flat, consistent with the insignificant mean spreads in Table~\ref{tab:portfolio_spreads}.}
\label{fig:cumulative}
\end{figure}

\FloatBarrier

%%------------------------------------------------------------
\subsection{Fama--MacBeth Cross-Sectional Regressions}
\label{sec:fm}
%%------------------------------------------------------------

To assess whether the portfolio-sort patterns survive standard risk-factor controls, we
estimate weekly Fama--MacBeth cross-sectional regressions. All slope coefficients are
averaged over time and $t$-statistics are constructed with Newey--West standard errors
(six lags). Controls comprise log market equity ($\log$ME), book-to-market (BM),
12-2 month momentum (MOM), short-term reversal (REV), idiosyncratic volatility (IVOL),
and Amihud (2002) illiquidity (ILLIQ), all winsorized at the 1st and 99th cross-sectional
percentiles each week. Table~\ref{tab:fm_aggregate} presents the baseline results.

\begin{table}[!htbp]
\centering
\caption{Fama--MacBeth Baseline Cross-Sectional Regressions (B1--B3)}
\label{tab:fm_aggregate}
\small
\begin{threeparttable}
\begin{tabular}{l c c c}
\toprule
 & B1 & B2 & B3 \\
\midrule
RSJ
  & \makecell{$-76.09^{***}$\\[-3pt]{\scriptsize$(-8.55)$}}
  & ---
  & \makecell{$-69.90^{***}$\\[-3pt]{\scriptsize$(-6.53)$}} \\[4pt]
RES
  & ---
  & \makecell{$-34.40$\\[-3pt]{\scriptsize$(-1.57)$}}
  & \makecell{$-54.53^{**}$\\[-3pt]{\scriptsize$(-2.48)$}} \\[5pt]
Controls (6 variables) & Yes & Yes & Yes \\[3pt]
Adj.\ $R^{2}$ (\%)     & 5.37 & 6.42 & 6.65 \\
$N_{\text{weeks}}$     & 1{,}590 & 1{,}462 & 1{,}462 \\
\bottomrule
\end{tabular}
\begin{tablenotes}
\footnotesize
\item \textit{Notes:} Time-averaged Fama--MacBeth slope coefficients (bps/week) with
Newey--West $t$-statistics (six lags) in parentheses. Specification B1 includes RSJ alone
with controls; B2 includes RES alone with controls; B3 includes both jointly. Controls
comprise $\log$(ME), BM, MOM, REV, IVOL, and log-Amihud illiquidity, all winsorized at
the 1st/99th cross-sectional percentiles weekly. The idiosyncratic volatility (IVOL)
coefficient in B1 is $-647.49$ bps ($t=-4.51$), replicating the \citet{ang2006} anomaly.
Control coefficients are suppressed for brevity; full coefficient vectors are available
upon request. Sample: January 1993--December 2024. \\
$^{*}p<0.10$,\ $^{**}p<0.05$,\ $^{***}p<0.01$.
\end{tablenotes}
\end{threeparttable}
\end{table}

Specification B1 confirms that RSJ earns a large and precisely estimated average slope of
$-76.09$ basis points ($t=-8.55$), explaining 5.37\% of weekly cross-sectional return
dispersion. The premium is not absorbed by any standard control: size, value, momentum,
reversal, idiosyncratic volatility, or liquidity. Notably, IVOL enters with a large
negative coefficient ($-647.49$ bps, $t=-4.51$), replicating the idiosyncratic volatility
anomaly documented by \citet{ang2006}. RSJ is distinct from these predictors as it
captures the directional sign asymmetry of intraday volatility rather than the level of
risk or recent price history.

Specification B2 shows that RES alone is statistically insignificant ($-34.40$ bps,
$t=-1.57$), even though its adjusted $R^2$ of 6.42\% is higher than that of RSJ. This
apparent contradiction resolves in Specification B3: including both RSJ and RES jointly
causes RES to nearly double in magnitude and reach statistical significance ($-54.53$
bps, $t=-2.48$), while RSJ remains highly significant ($-69.90$ bps, $t=-6.53$). When
estimated alone, RES's predictive content is masked because it co-moves with controls
--- especially IVOL and size --- that are also correlated with the omitted RSJ signal.
The joint adjusted $R^2$ of 6.65\% is larger than either univariate specification,
indicating that tail severity and volatility asymmetry capture distinct dimensions of
downside risk.

Figure~\ref{fig:ts_betas} plots the rolling 52-week average of the weekly FM slopes for
RSJ and RES from specification B3. Two features stand out. First, the RSJ slope is
persistently negative across the entire 32-year sample, never crossing zero in any
sustained period. This constancy distinguishes RSJ from most other return predictors,
whose cross-sectional pricing tends to be highly regime-dependent. Second, the RES slope
is markedly volatile: it spikes sharply positive during the 2000--2002 technology bust
and the 2008--2009 financial crisis before reverting to negative values during
expansions. This temporal pattern anticipates the crisis-state analysis in
Section~\ref{sec:crisis} and suggests that RSJ and RES pricing operate through
mechanisms with different degrees of time variation.

\begin{figure}[!ht]
\centering
\includegraphics[width=0.88\textwidth]{results/figures/fig_time_series_fm_betas_b3}
\caption{Rolling 52-week mean of weekly Fama--MacBeth slope coefficients (bps) for RSJ
(blue) and RES (orange) from the joint specification B3 (RSJ $+$ RES $+$ controls).
Shaded vertical bands mark NBER recession periods. The RSJ slope is persistently negative
across the full sample. The RES slope is highly time-varying, spiking sharply positive
during the 2000--2002 and 2008--2009 downturns and reverting to negative during
expansions, consistent with the crisis-state sign reversal in Table~\ref{tab:crisis_state}.}
\label{fig:ts_betas}
\end{figure}

\FloatBarrier

%%------------------------------------------------------------
\subsection{Decomposition Analysis and Formal Tests}
\label{sec:decomp}
%%------------------------------------------------------------

Having established the baseline premia, we now decompose RSJ and RES into their systematic
and idiosyncratic components to identify the source of the predictability.
Table~\ref{tab:fm_decomposition} presents all five decomposition specifications: B4 and B5
test the RSJ decomposition (M1 and M2) in isolation; B6 tests the RES decomposition alone;
B7 (Method~1) and B8 (Method~2) are the primary empirical specifications, simultaneously
including all four decomposed components and the full control set.

\begin{table}[!htbp]
\centering
\caption{Fama--MacBeth Decomposition Regressions (B4--B8)}
\label{tab:fm_decomposition}
\small
\begin{threeparttable}
\begin{tabular}{l c c c c c}
\toprule
 & B4 & B5 & B6 & B7 & B8 \\
 & (M1: RSJ) & (M2: RSJ) & (RES) & (M1: Full) & (M2: Full) \\
\midrule
RSJ$^{\text{sys}}$ (M1)
  & \makecell{$-3.15$\\[-3pt]{\scriptsize$(-0.78)$}}
  & --- & ---
  & \makecell{$-3.15$\\[-3pt]{\scriptsize$(-0.78)$}}
  & --- \\[4pt]
RSJ$^{\text{idio}}$ (M1)
  & \makecell{$-43.70^{***}$\\[-3pt]{\scriptsize$(-4.44)$}}
  & --- & ---
  & \makecell{$-43.70^{***}$\\[-3pt]{\scriptsize$(-4.44)$}}
  & --- \\[4pt]
RSJ$^{\text{sys}}$ (M2)
  & --- & \makecell{$+17.31$\\[-3pt]{\scriptsize$(+0.46)$}} & ---
  & --- & \makecell{$+17.31$\\[-3pt]{\scriptsize$(+0.46)$}} \\[4pt]
RSJ$^{\text{idio}}$ (M2)
  & --- & \makecell{$-58.71^{***}$\\[-3pt]{\scriptsize$(-5.78)$}} & ---
  & --- & \makecell{$-58.71^{***}$\\[-3pt]{\scriptsize$(-5.78)$}} \\[4pt]
RES$^{\text{sys}}$
  & --- & ---
  & \makecell{$---^{b}$\\[-3pt]{\scriptsize$(-0.66)$}}
  & \makecell{$---^{b}$\\[-3pt]{\scriptsize$(-0.66)$}}
  & \makecell{$---^{b}$\\[-3pt]{\scriptsize$(-0.47)$}} \\[4pt]
RES$^{\text{idio}}$
  & --- & ---
  & \makecell{$-56.67^{***}$\\[-3pt]{\scriptsize$(-3.05)$}}
  & \makecell{$-56.67^{***}$\\[-3pt]{\scriptsize$(-3.05)$}}
  & \makecell{$-66.32^{***}$\\[-3pt]{\scriptsize$(-3.41)$}} \\[5pt]
Controls (6)        & Yes & Yes & Yes & Yes & Yes \\[2pt]
Adj.\ $R^{2}$ (\%) & ---  & ---  & ---  & 8.38 & 8.45 \\
$N_{\text{weeks}}$ & 1{,}463 & 1{,}411 & 1{,}462 & 1{,}463 & 1{,}411 \\
\bottomrule
\end{tabular}
\begin{tablenotes}
\footnotesize
\item \textit{Notes:} Time-averaged Fama--MacBeth slope coefficients (bps/week) with
Newey--West $t$-statistics (six lags) in parentheses. B4 and B5 include only the
systematic and idiosyncratic RSJ components (M1 and M2, respectively) with controls; B6
includes only the systematic and idiosyncratic RES components; B7 and B8 are the primary
full decomposition specifications including all four components simultaneously. Controls
are identical to those in Table~\ref{tab:fm_aggregate}. ``---'' in a predictor cell
indicates that the variable is not included in that specification; Adj.~$R^{2}$ is not
reported for the partial decompositions (B4--B6) as these are intermediate robustness
checks. \\
$^{b}$ Systematic RES is estimated with large standard errors due to collinearity with
aggregate market conditions and the control set; the coefficient is therefore not
numerically interpreted. Only the $t$-statistic is reported. \\
$^{*}p<0.10$,\ $^{**}p<0.05$,\ $^{***}p<0.01$.
\end{tablenotes}
\end{threeparttable}
\end{table}

The results are unambiguous. For RSJ, the idiosyncratic component is highly significant
under both decomposition methods (M1: $-43.70$ bps, $t=-4.44$; M2: $-58.71$ bps,
$t=-5.78$). Systematic RSJ, in contrast, is statistically indistinguishable from zero
and even reverses sign across methods (M1: $-3.15$, $t=-0.78$; M2: $+17.31$,
$t=+0.46$). The same pattern holds for RES: idiosyncratic RES is robustly significant
(M1: $-56.67$ bps, $t=-3.05$; M2: $-66.32$ bps, $t=-3.41$), while systematic RES,
though large in magnitude, is estimated with too much imprecision to be distinguished
from zero under either method ($t=-0.66$ and $t=-0.47$, respectively). The wide
confidence intervals on the systematic components reflect their high collinearity with
aggregate market conditions and the controls already in the regression.

The adjusted $R^2$ of 8.38\% (M1) and 8.45\% (M2) substantially exceeds the 6.65\%
achieved by the joint aggregate specification B3. This improvement of approximately
1.8 percentage points confirms that separating the measures into systematic and
idiosyncratic dimensions recovers additional cross-sectional information that is
subsumed when aggregate measures are used. Critically, this gain is attributable to the
idiosyncratic components: they isolate firm-specific variation that accounts for
virtually all predictive content, allowing the signal to be extracted more cleanly.

These patterns are illustrated in Figure~\ref{fig:fm_coef}, which plots the point
estimates and 95\% confidence intervals from specifications B3, B7, and B8. The
idiosyncratic components cluster tightly to the left of zero; the systematic components
span zero with substantially wider bands.

\begin{figure}[!ht]
\centering
\includegraphics[width=0.80\textwidth]{results/figures/fig_fm_coefficients_selected}
\caption{Fama--MacBeth slope coefficients (bps/week) with 95\% Newey--West confidence
intervals from the aggregate joint specification (B3) and the full decomposition
specifications (B7 and B8). Idiosyncratic components are tightly estimated to the left
of zero. Systematic components span zero with substantially wider bands, consistent with
the Wald test results in Table~\ref{tab:wald_tests}.}
\label{fig:fm_coef}
\end{figure}

We formalise the comparison between systematic and idiosyncratic risk prices using Wald
tests on the time-averaged coefficient vector from B7 and B8. The four hypotheses follow
directly from the methodology: $H_1$ tests whether systematic and idiosyncratic RSJ
carry equal risk prices; $H_2$ tests the same for RES; $H_3$ tests the joint significance
of the idiosyncratic block; and $H_4$ tests the joint significance of the systematic
block. Table~\ref{tab:wald_tests} reports the Wald statistics.

\begin{table}[!htbp]
\centering
\caption{Wald Tests: Systematic versus Idiosyncratic Downside Risk}
\label{tab:wald_tests}
\small
\begin{threeparttable}
\begin{tabular}{l l c c c c}
\toprule
 & Null hypothesis & df
   & $\chi^2$ (M1) & $\chi^2$ (M2) & Decision \\
\midrule
$H_1$ & $\bar\beta^{\text{sys}}_{\text{RSJ}} = \bar\beta^{\text{idio}}_{\text{RSJ}}$
  & 1 & $16.65^{***}$ & $4.77^{**}$ & Reject \\[3pt]
$H_2$ & $\bar\beta^{\text{sys}}_{\text{RES}} = \bar\beta^{\text{idio}}_{\text{RES}}$
  & 1 & $0.06$ & $0.00$ & Fail to reject \\[3pt]
$H_3$ & $\bar\beta^{\text{idio}}_{\text{RSJ}} = \bar\beta^{\text{idio}}_{\text{RES}} = 0$
  & 2 & $28.34^{***}$ & $41.48^{***}$ & Reject \\[3pt]
$H_4$ & $\bar\beta^{\text{sys}}_{\text{RSJ}} = \bar\beta^{\text{sys}}_{\text{RES}} = 0$
  & 2 & $1.08$ & $0.34$ & Fail to reject \\
\bottomrule
\end{tabular}
\begin{tablenotes}
\footnotesize
\item \textit{Notes:} Wald statistics computed from the time-averaged coefficient
vector $\bar{\bm{\beta}}$ and its Newey--West (six lags) estimated covariance matrix
from the primary full decomposition specifications B7 (M1) and B8 (M2). The statistic
is asymptotically $\chi^2(q)$ under the null, where $q$ is the number of
restrictions. $H_1$ and $H_2$ test equality of risk prices between systematic and
idiosyncratic components within RSJ and RES, respectively. $H_3$ jointly tests whether
both idiosyncratic components are unpriced; $H_4$ jointly tests whether both systematic
components are unpriced. \\
$^{*}p<0.10$,\ $^{**}p<0.05$,\ $^{***}p<0.01$.
\end{tablenotes}
\end{threeparttable}
\end{table}

The contrast between $H_3$ and $H_4$ constitutes the core empirical finding of this
study. $H_3$ --- that both idiosyncratic components are jointly zero --- is decisively
rejected under both methods ($\chi^2 = 28.34^{***}$ under M1; $41.48^{***}$ under M2).
$H_4$ --- that both systematic components are jointly zero --- cannot be rejected at any
conventional significance level ($\chi^2 = 1.08$ and $0.34$). The predictive content of
asymmetric volatility and tail risk in the cross-section of stock returns is therefore
located entirely in firm-specific variation.

For RSJ, $H_1$ is rejected at the 1\% level under M1 ($\chi^2 = 16.65$) and the 5\%
level under M2 ($\chi^2 = 4.77$), confirming that systematic and idiosyncratic RSJ are
priced significantly differently. For RES, $H_2$ cannot be rejected under either method
($\chi^2 = 0.06$ and $0.00$). This does not undermine the idiosyncratic interpretation:
both RES components carry negative risk prices, and the failure to reject $H_2$
reflects insufficient precision in the systematic estimate rather than equivalence of
the underlying channels. The significant $H_3$ confirms that idiosyncratic tail risk is
robustly priced; the failure of $H_4$ confirms that systematic tail risk is not.

Supplementary specifications B9 and B10 in Table~\ref{tab:fm_supplementary} of
Appendix~\ref{app:tables} further solidify this conclusion: when idiosyncratic RSJ
and idiosyncratic RES are estimated jointly (B9), both remain significant ($-41.74$
bps, $t=-4.27$ and $-54.71$ bps, $t=-2.64$), whereas replacing idiosyncratic RSJ with
its systematic counterpart (B10) renders that component negligible ($-0.94$ bps,
$t=-0.22$).

The temporal stability of the idiosyncratic premia is examined in
Figure~\ref{fig:rolling_idio}. The rolling 52-week FM slopes for idiosyncratic RSJ (M1)
and idiosyncratic RES from specification B7 are both persistently negative throughout
the sample, with no structural break in 2013 when the \citet{bollerslev2020} sample
ends. This out-of-sample persistence is a demanding robustness test: the same
decomposition protocol applied to an additional eleven years of data yields premia of
comparable magnitude to those in the original window. The finding rules out the
possibility that the idiosyncratic channel is a post-hoc artefact specific to the
training sample used to motivate the RSJ measure.

\begin{figure}[!ht]
\centering
\includegraphics[width=0.88\textwidth]{results/figures/fig_rolling_fm_decomp_idio}
\caption{Rolling 52-week mean of weekly Fama--MacBeth slope coefficients (bps) for
idiosyncratic RSJ (M1, blue) and idiosyncratic RES (orange) from the full decomposition
specification B7. Shaded bands mark NBER recessions. Both idiosyncratic premia are
persistently negative throughout the sample, including the post-2013 period not covered
by \citet{bollerslev2020}, confirming that the idiosyncratic channel is not an artefact
of the original estimation window.}
\label{fig:rolling_idio}
\end{figure}

\FloatBarrier

%%------------------------------------------------------------
\subsection{Crisis-State Analysis}
\label{sec:crisis}
%%------------------------------------------------------------

A risk-based interpretation predicts that systematic downside risk measures should earn
higher premia during aggregate downturns, when investors' marginal utility is elevated.
We test this by interacting the systematic components of RSJ and RES with a dummy for
NBER recession weeks. The sample spans three recession episodes: March--November 2001,
December 2007--June 2009, and February--April 2020, totalling 134 recession weeks.

Table~\ref{tab:crisis_state} reports the primary crisis-state specification, which
jointly includes systematic RSJ (M1) and systematic RES with NBER interactions. The
columns report the time-averaged FM slopes for non-recession and recession weeks
separately, along with the crisis increment and its $t$-statistic.

\begin{table}[!htbp]
\centering
\caption{Crisis-State Fama--MacBeth Regressions: Systematic Components}
\label{tab:crisis_state}
\small
\begin{threeparttable}
\begin{tabular}{l c c c c}
\toprule
Predictor
  & Expansion (bps) & Recession (bps)
  & $\Delta$ (bps) & $t(\Delta)$ \\
\midrule
RSJ$^{\text{sys}}$ (M1) & $-2.07$ & $+8.44$  & $+10.51$ & $(0.46)$ \\[3pt]
RES$^{\text{sys}}$      & $-91.46$ & $-234.03$ & $-142.57$ & $(-0.29)$ \\
\bottomrule
\end{tabular}
\begin{tablenotes}
\footnotesize
\item \textit{Notes:} Time-averaged Fama--MacBeth slope coefficients (bps/week) from
the joint crisis interaction specification (equation~\eqref{eq:crisis}), which includes
systematic RSJ (M1), systematic RES, an NBER recession dummy, and the interactions of
each systematic predictor with the recession dummy, alongside the full control set.
Expansion reports the average slope during non-recession weeks; Recession reports the
average slope during recession weeks; $\Delta$ is the crisis-minus-expansion increment.
$t(\Delta)$ uses Newey--West standard errors (six lags). NBER recession periods:
March--November 2001, December 2007--June 2009, and February--April 2020 (134 weeks).
The complete set of crisis-state results across all specifications and decomposition
methods is in Table~\ref{tab:crisis_full} of Appendix~\ref{app:tables}. \\
$^{*}p<0.10$,\ $^{**}p<0.05$,\ $^{***}p<0.01$.
\end{tablenotes}
\end{threeparttable}
\end{table}

The results provide no evidence of crisis-state amplification for either systematic
measure. The systematic RSJ slope shifts from $-2.07$ bps in expansions to $+8.44$ bps
in recessions --- a sign reversal rather than an intensification --- but the crisis
increment of $+10.51$ bps is statistically negligible ($t=0.46$). Systematic RES shows
no meaningful change either: the slope is $-91.46$ bps in expansions and $-234.03$ bps
in recessions, but the crisis increment of $-142.57$ bps carries a $t$-statistic of
$-0.29$. Neither systematic component intensifies during recessions, which represents
the most direct challenge to a rational risk-premium explanation: compensation for
aggregate downside exposure should be most valuable precisely when aggregate conditions
deteriorate. Under any stochastic discount factor framework in which aggregate bad states carry a
positive risk price, the premium on systematic downside exposure should rise when that
exposure is most painful to bear. The absence of any such amplification across three
distinct recession episodes spanning very different economic environments---a
technology-sector contraction, a global banking crisis, and a pandemic shock---constitutes
a broad-based failure of this prediction.

The complete crisis-state analysis across all predictor specifications and decomposition
methods is reported in Table~\ref{tab:crisis_full} of Appendix~\ref{app:tables}. One
additional pattern in those results is notable: the idiosyncratic component of RES
undergoes a statistically significant sign reversal during recessions. In the full
decomposition specifications, the idiosyncratic RES slope is strongly negative during
expansions ($-66.94$ to $-78.49$ bps) but turns positive during recessions, with crisis
increments of $+112.04$ bps ($t=2.01$, M1) and $+128.13$ bps ($t=2.28$, M2). This
reversal is consistent with a mispricing correction: stocks that are systematically
overpriced in normal times --- and therefore earn negative abnormal returns --- recover
when tail risks actually materialise in aggregate downturns, as the mispricing unwinds
\citep{atilgan2020}.

\FloatBarrier

%%------------------------------------------------------------
\subsection{Limits to Arbitrage and Market Frictions}
\label{sec:frictions}
%%------------------------------------------------------------

If idiosyncratic downside risk is mispriced, its persistence should depend on the cost
of correcting it. The limits-to-arbitrage hypothesis predicts that predictability
concentrates among stocks where trading frictions are most severe \citep{shleifer1997}.
Two tests examine this prediction.

\paragraph{Idiosyncratic characteristics and illiquidity.}
As a first step, we assess whether idiosyncratic risk measures are systematically higher
for illiquid stocks. Table~\ref{tab:illiquidity} reports cross-sectional regressions of
each idiosyncratic signal on log-Amihud illiquidity, with and without firm-level controls.

\begin{table}[!htbp]
\centering
\caption{Idiosyncratic Downside Risk and Illiquidity}
\label{tab:illiquidity}
\small
\begin{threeparttable}
\begin{tabular}{l c c c c c}
\toprule
Dependent variable
  & \multicolumn{2}{c}{Univariate}
  & \multicolumn{2}{c}{Controlled}
  & $N_{\text{weeks}}$ \\
\cmidrule(lr){2-3}\cmidrule(lr){4-5}
  & $\hat\delta$ & $t$ & $\hat\delta$ & $t$ & \\
\midrule
RSJ$^{\text{idio}}$ (M1)
  & $-0.0005^{*}$ & $(-2.36)$ & $0.0042^{***}$ & $(9.22)$ & 1{,}464 \\[3pt]
RES$^{\text{idio}}$
  & $0.0378^{***}$ & $(22.67)$ & $0.0163^{***}$ & $(28.92)$ & 1{,}463 \\
\bottomrule
\end{tabular}
\begin{tablenotes}
\footnotesize
\item \textit{Notes:} Time-averaged slope coefficients from weekly cross-sectional
regressions of each idiosyncratic downside-risk signal on log-Amihud illiquidity
(equation~\eqref{eq:mech}). The univariate specification includes only log-Amihud
illiquidity as the regressor. The controlled specification adds $\log$(ME), BM, MOM,
REV, and IVOL; illiquidity is excluded from the control set to avoid double-counting.
$t$-statistics use Newey--West standard errors (six lags). The sign reversal for
RSJ$^{\text{idio}}$ (M1)---negative univariate, positive controlled---reflects that
the raw negative correlation is driven by offsetting size and volatility effects;
conditional on firm characteristics, more illiquid stocks exhibit greater idiosyncratic
downside asymmetry. Results for RSJ$^{\text{idio}}$ (M2) are reported in
Table~\ref{tab:frictions_signal} of Appendix~\ref{app:tables}. \\
$^{*}p<0.10$,\ $^{**}p<0.05$,\ $^{***}p<0.01$.
\end{tablenotes}
\end{threeparttable}
\end{table}

Idiosyncratic RES is strongly and positively associated with illiquidity in both
specifications: the univariate slope is $0.0378$ ($t=22.67$), and it remains $0.0163$
($t=28.92$) after controlling for firm characteristics. Illiquid stocks systematically
exhibit greater idiosyncratic tail risk, consistent with a setting in which limited
institutional attention or information barriers concentrate firm-specific negative events
in harder-to-trade names. For idiosyncratic RSJ (M1), the univariate association is
weakly negative ($-0.0005$, $t=-2.36$), but it reverses to a positive and highly
significant controlled estimate ($+0.0042$, $t=9.22$). This sign reversal indicates
that the raw negative correlation reflects offsetting effects of size and volatility:
conditional on firm characteristics, more illiquid stocks exhibit greater idiosyncratic
downside asymmetry.

\paragraph{Illiquidity interaction tests.}
The second test directly examines whether the return-predictive content of RSJ is
amplified among illiquid stocks. Table~\ref{tab:frictions} reports Fama--MacBeth
interaction regressions of the form $R_{i,w+1} = \alpha + \beta_Z Z_{i,w}
+ \beta_\text{I}\,\text{ILLIQ}_{i,w} + \beta_\times (Z_{i,w} \times \text{ILLIQ}_{i,w})
+ \gamma' X_{i,w} + \varepsilon_{i,w+1}$, where controls $X$ include $\log(\text{ME})$,
BM, MOM, REV, and IVOL. The interaction coefficient $\beta_\times$ measures whether the
return-predictive slope of signal $Z$ varies systematically with trading frictions.

\begin{table}[!htbp]
\centering
\begin{threeparttable}
\caption{Market Frictions: Fama--MacBeth Interaction Regressions}
\label{tab:frictions}
\small
\begin{tabular}{l c c c r}
\toprule
Signal
  & $\hat\beta_Z$ (bps) & $\hat\beta_{\text{ILLIQ}}$ (bps)
  & $\hat\beta_\times$ (bps) & $N_{\text{wks}}$ \\
\midrule
\multicolumn{5}{l}{\textit{Panel A: RSJ signals}} \\[2pt]
RSJ Total
  & \makecell{$-444.55^{***}$\\[-3pt]{\scriptsize$(-7.94)$}}
  & \makecell{$-7.56^{***}$\\[-3pt]{\scriptsize$(-5.77)$}}
  & \makecell{$-18.73^{***}$\\[-3pt]{\scriptsize$(-6.86)$}}
  & 1{,}590 \\[4pt]
RSJ Idiosyncratic (M1)
  & \makecell{$-314.55^{***}$\\[-3pt]{\scriptsize$(-4.00)$}}
  & \makecell{$-6.96^{***}$\\[-3pt]{\scriptsize$(-4.83)$}}
  & \makecell{$-13.48^{***}$\\[-3pt]{\scriptsize$(-3.65)$}}
  & 1{,}463 \\[4pt]
RSJ Idiosyncratic (M2)
  & \makecell{$-308.34^{***}$\\[-3pt]{\scriptsize$(-6.09)$}}
  & \makecell{$-6.93^{***}$\\[-3pt]{\scriptsize$(-4.91)$}}
  & \makecell{$-12.74^{***}$\\[-3pt]{\scriptsize$(-5.08)$}}
  & 1{,}411 \\[5pt]
\multicolumn{5}{l}{\textit{Panel B: RES signals}} \\[2pt]
RES Total
  & \makecell{$-95.09$\\[-3pt]{\scriptsize$(-0.74)$}}
  & \makecell{$-6.41^{***}$\\[-3pt]{\scriptsize$(-3.57)$}}
  & \makecell{$-3.96$\\[-3pt]{\scriptsize$(-0.56)$}}
  & 1{,}462 \\[4pt]
RES Idiosyncratic
  & \makecell{$-82.24$\\[-3pt]{\scriptsize$(-0.65)$}}
  & \makecell{$-6.58^{***}$\\[-3pt]{\scriptsize$(-3.69)$}}
  & \makecell{$-3.07$\\[-3pt]{\scriptsize$(-0.44)$}}
  & 1{,}462 \\
\bottomrule
\end{tabular}
\begin{tablenotes}
\footnotesize
\item \textit{Notes:} Each row reports a separate Fama--MacBeth regression including
the signal $Z$, log-Amihud illiquidity (ILLIQ), the interaction $Z \times \text{ILLIQ}$,
and firm-level controls $\log(\text{ME})$, BM, MOM, REV, IVOL. All coefficients in
bps/week; $t$-statistics in parentheses use Newey--West standard errors (six lags).
The full table including systematic components is in Table~\ref{tab:frictions_full}
of Appendix~\ref{app:tables}.
$^{*}p<0.10$,\ $^{**}p<0.05$,\ $^{***}p<0.01$.
\end{tablenotes}
\end{threeparttable}
\end{table}

For RSJ, both the main effect and the illiquidity interaction are large, negative, and
highly significant. The interaction coefficient for RSJ total is $-18.73$ bps
($t=-6.86$), indicating that the return-predictive slope of RSJ is meaningfully stronger
for less liquid stocks. For idiosyncratic RSJ, this amplification is similar in magnitude
and robust under both decompositions (M1: $-13.48$ bps, $t=-3.65$; M2: $-12.74$ bps,
$t=-5.08$). Because log-Amihud illiquidity spans roughly four standard deviations in the
cross-section, moving a stock from the median to the 75th illiquidity percentile amplifies
the idiosyncratic RSJ slope by approximately $6$--$9$ basis points per unit of RSJ ---
an economically meaningful amplification. These results confirm that the idiosyncratic RSJ
premium is largest precisely where arbitrage is most costly, consistent with the
limits-to-arbitrage interpretation.

For RES, the interaction coefficients are economically small and statistically
insignificant for both the total and idiosyncratic variants ($-3.96$, $t=-0.56$; $-3.07$,
$t=-0.44$). The return-predictive content of RES does not vary systematically with
illiquidity in the interaction sense, even though the level of idiosyncratic RES itself
co-moves strongly with illiquidity (Table~\ref{tab:illiquidity}). The asymmetry between
RSJ and RES in the friction test suggests that the two measures capture distinct aspects
of the mispricing channel. For RSJ, the mispricing channel operates primarily through
the \textit{predictive slope}: less liquid stocks have a steeper RSJ--return gradient,
meaning the same degree of downside asymmetry generates a larger subsequent return
differential when arbitrage is costly. For RES, the friction mechanism operates instead
through \textit{exposure}: illiquid stocks are systematically more exposed to
idiosyncratic tail risk, but conditional on that level of exposure, the return
predictability of RES does not further depend on how hard it is to trade the stock.
This distinction suggests that RSJ and RES embed information about mispricing through
different economic pathways, and that a complete account of firm-specific downside risk
pricing requires both measures rather than either one alone.

Taken together, the evidence from all five analyses converges on a coherent
interpretation. Idiosyncratic asymmetric volatility and idiosyncratic tail risk are
robustly priced in the cross-section; their systematic counterparts are not. No
crisis-state amplification of the systematic premia is detectable, contrary to the
predictions of a rational risk-compensation view. The idiosyncratic RSJ premium is
amplified among illiquid stocks, and idiosyncratic RES co-moves with trading frictions.
These findings are mutually reinforcing: firm-specific downside shocks generate
mispricings that persist precisely because the stocks bearing the most idiosyncratic
tail risk are those where arbitrage is most costly.

\FloatBarrier

%%============================================================
%%  CONCLUSION
%%============================================================
\section{Conclusion}
\label{sec:conclusion}

This paper investigates the mechanism driving the return predictability of realized
asymmetric volatility and contributes two innovations to the literature. First, it
introduces realized expected shortfall, constructed from high-frequency intraday returns
under the self-similar scaling framework of \citet{Dimitriadis03072022}, as a
complementary measure of left-tail severity. Second, it decomposes both RSJ and RES into
systematic and idiosyncratic components using an intraday market model, enabling the
first structured empirical assessment of whether the RSJ anomaly reflects rational risk
compensation or limits-to-arbitrage-induced mispricing.

The evidence from five separate analyses --- portfolio sorts, Fama--MacBeth regressions,
Wald tests, crisis-state regressions, and friction interaction tests --- converges on a
coherent picture. The idiosyncratic channel dominates throughout. In quintile portfolio
sorts, idiosyncratic RSJ earns 17--20 basis points per week under both Method~1 and
Method~2, essentially matching the total RSJ spread, while systematic RSJ is
statistically indistinguishable from zero under every specification. The same pattern
holds in Fama--MacBeth regressions, where idiosyncratic RSJ commands slopes of $-40$ to
$-59$ basis points per week compared with a near-zero, sign-unstable systematic slope.
Separating the decomposed components raises the average adjusted $R^2$ from 6.65\%
in the joint aggregate specification to 8.38--8.45\% in the full decomposition models,
confirming that the decomposition is economically consequential and not merely a
reparameterisation.

Wald tests formalise this asymmetry. The joint null that both idiosyncratic components
are unpriced is rejected decisively under both decomposition methods ($\chi^2 = 28.34$
and $41.48$, both significant at the 1\% level). The joint null that both systematic
components are unpriced cannot be rejected ($\chi^2 = 1.08$ and $0.34$), implying that
the systematic channel contributes nothing beyond what the controls already capture.
Further, the equality of systematic and idiosyncratic slopes within RSJ is formally
rejected ($\chi^2 = 16.65^{***}$ under M1), whereas the same hypothesis for RES cannot
be rejected, revealing a striking asymmetry between the two measures.

The crisis-state evidence is the most direct challenge to a risk-premium view. If RSJ
reflects compensation for systematic downside exposure, its premium should be amplified
precisely when aggregate conditions are most adverse --- during recessions, when bearing
such exposure is costliest. We find no evidence of this. The interaction between
systematic RSJ and the NBER recession indicator is economically small and statistically
insignificant ($t = 0.46$), and the same holds for systematic RES ($t = -0.29$). The
predictive content of the systematic channel does not respond to the business cycle in
any way consistent with rational risk compensation. If anything, the idiosyncratic RES
premium shows a sign reversal during recessions, consistent with a correction of
accumulated tail-risk mispricing when tail risks actually materialise.

The frictions tests complete the picture by establishing a mechanism. The idiosyncratic
RSJ premium is significantly stronger among illiquid stocks: the interaction with
log-Amihud illiquidity is $-13.48$ basis points per unit of RSJ ($t = -3.65$) under
Method~1 and $-12.74$ ($t = -5.08$) under Method~2. Moving a stock from the median to
the 75th illiquidity percentile amplifies the idiosyncratic RSJ slope by approximately
6--9 basis points per unit of signal, an amplification that is both statistically and
economically meaningful. Idiosyncratic RES, by contrast, is itself strongly associated
with illiquidity --- illiquid stocks exhibit more severe intraday tail risk --- but its
return-predictive content does not interact with frictions in the same way, suggesting
that the two measures operate through partially distinct aspects of the mispricing
channel. Together, these findings indicate that firm-specific downside shocks generate
persistent mispricings because the stocks most exposed to idiosyncratic tail risk are
precisely those where the cost of arbitrage is highest.

Several limitations bound the scope of these conclusions. The intraday market model uses
SPY as a market proxy, and the quality of this decomposition may vary across market
regimes and firm-size quintiles. The analysis is restricted to U.S. equities; whether
the same systematic-versus-idiosyncratic pattern holds in international markets, where
institutional structures and arbitrage conditions differ, remains an open question. The
identification of the mechanism as mispricing is inferential: we observe the
cross-sectional and time-series patterns that a mispricing view predicts, but we do not
directly identify the cognitive or institutional frictions that generate the initial
underreaction to firm-specific tail risk.

Future work could exploit exogenous variation in arbitrage costs --- such as short-sale
constraint changes or the entry and exit of institutional investors --- to sharpen the
causal identification. Extending the RES measure to international equity panels and
testing whether it contains incremental information beyond existing tail risk measures in
bond and derivative markets are further natural extensions. The question of whether the
idiosyncratic RSJ premium is driven by investor inattention, analyst neglect, or
systematic biases in probability assessment for extreme outcomes represents a promising
avenue for integrating the results of this paper with the broader behavioural finance
literature.

%%============================================================
%%  APPENDIX
%%============================================================
\appendix
\clearpage
\section{Appendix A: Additional Figures}
\label{app:figures}

\begin{figure}[!ht]
\centering
\includegraphics[width=0.78\textwidth]{results/figures/fig_predictor_correlation_heatmap}
\caption{Pairwise Pearson correlations among all return predictors in the weekly panel.
RSJ and its idiosyncratic component are highly correlated; the systematic component is
weakly correlated with both. RES displays a similar within-group pattern. Cross-group
correlations between RSJ and RES measures are low, supporting their treatment as
distinct predictors in the joint specifications.}
\label{fig:corr_heatmap}
\end{figure}

\begin{figure}[!ht]
\centering
\includegraphics[width=0.88\textwidth]{results/figures/fig_weekly_cross_section_coverage}
\caption{Number of stocks in the estimation panel by week, January 1993 to December
2024. Coverage grows from approximately 1{,}000 stocks in the early 1990s to a peak of
roughly 3{,}200 around 2000, declines gradually through the 2000s and 2010s, and
stabilises at approximately 2{,}000--2{,}500 in recent years, reflecting the secular
decline in the number of listed U.S.\ equities after the late 1990s.}
\label{fig:coverage}
\end{figure}

%%------------------------------------------------------------
\clearpage
\section{Appendix B: Additional Tables}
\label{app:tables}
%%------------------------------------------------------------

\begin{table}[!htbp]
\centering
\caption{Summary Statistics: Full Sample (January 1993 -- December 2024)}
\label{tab:sumstats_full}
\small
\begin{threeparttable}
\begin{tabular}{l r r r r r r}
\toprule
Variable & $N$ & Mean & Std & P25 & Median & P75 \\
\midrule
\multicolumn{7}{l}{\textit{Panel A: Returns}} \\[2pt]
$R_{i,w}$
  & 3{,}799{,}481 & $0.0060$ & $0.0997$ & $-0.0281$ & $0.0025$ & $0.0341$ \\
$R_{i,w+1}$
  & 3{,}799{,}481 & $0.0024$ & $0.0759$ & $-0.0290$ & $0.0013$ & $0.0317$ \\[5pt]
\multicolumn{7}{l}{\textit{Panel B: RSJ measures}} \\[2pt]
RSJ Total
  & 3{,}807{,}204 & $0.0047$ & $0.1388$ & $-0.0832$ & $0.0045$ & $0.0924$ \\
RSJ Systematic (M2)
  & 3{,}142{,}050 & $0.0043$ & $0.0476$ & $-0.0226$ & $0.0041$ & $0.0313$ \\
RSJ Idiosyncratic (M2)
  & 3{,}142{,}050 & $0.0006$ & $0.1343$ & $-0.0851$ & $0.0004$ & $0.0862$ \\
RSJ Systematic (M1)
  & 3{,}152{,}418 & $0.0198$ & $0.3568$ & $-0.1334$ & $0.0125$ & $0.1628$ \\
RSJ Idiosyncratic (M1)
  & 3{,}152{,}418 & $0.0025$ & $0.1431$ & $-0.0894$ & $0.0025$ & $0.0946$ \\[5pt]
\multicolumn{7}{l}{\textit{Panel C: RES measures}} \\[2pt]
RES Total
  & 3{,}167{,}885 & $0.1826$ & $0.1468$ & $0.0934$ & $0.1388$ & $0.2193$ \\
RES Systematic
  & 3{,}167{,}885 & $0.0357$ & $0.0328$ & $0.0156$ & $0.0272$ & $0.0451$ \\
RES Idiosyncratic
  & 3{,}167{,}885 & $0.1764$ & $0.1469$ & $0.0877$ & $0.1320$ & $0.2114$ \\[5pt]
\multicolumn{7}{l}{\textit{Panel D: Firm characteristics}} \\[2pt]
$\log(\text{ME})$
  & 3{,}807{,}204 & $14.146$ & $1.630$  & $12.962$ & $13.981$ & $15.144$ \\
BM
  & 3{,}435{,}338 & $0.553$  & $0.540$  & $0.249$  & $0.442$  & $0.717$  \\
MOM
  & 3{,}574{,}018 & $0.248$  & $0.927$  & $-0.125$ & $0.108$  & $0.385$  \\
REV
  & 3{,}801{,}129 & $0.0054$ & $0.0837$ & $-0.0278$& $0.0022$ & $0.0332$ \\
IVOL
  & 3{,}783{,}933 & $0.0223$ & $0.0205$ & $0.0112$ & $0.0173$ & $0.0275$ \\
ILLIQ (log)
  & 3{,}801{,}415 & $-20.114$& $2.068$  & $-21.533$& $-20.029$& $-18.610$\\[5pt]
\multicolumn{7}{l}{\textit{Panel E: Data quality}} \\[2pt]
Valid trading days / week
  & 3{,}807{,}204 & $4.761$ & $0.790$ & $5.000$ & $5.000$ & $5.000$ \\
Intraday obs.\ / stock-week
  & 3{,}807{,}204 & $6{,}036$ & $12{,}717$ & $724$ & $2{,}132$ & $6{,}199$ \\
\bottomrule
\end{tabular}
\begin{tablenotes}
\footnotesize
\item \textit{Notes:} Summary statistics for all stock-week observations in the full
1993--2024 panel. RSJ measures the relative signed jump as the normalised difference
between positive and negative realised semivariance. M2 isolates the systematic component
via a rolling 52-week market-RSJ regression; M1 decomposes signed intraday jumps prior
to constructing RSJ. RES is weekly realised expected shortfall ($p=0.025$, $H=0.5$,
sign-adjusted). All controls are winsorized at the 1st/99th percentiles weekly.
The Bollerslev-comparable 1993--2013 subsample is in Table~\ref{tab:sumstats_bollerslev}.
\end{tablenotes}
\end{threeparttable}
\end{table}

\begin{table}[!htbp]
\centering
\caption{Summary Statistics: Bollerslev-Comparable Subsample (1993 -- 2013)}
\label{tab:sumstats_bollerslev}
\small
\begin{threeparttable}
\begin{tabular}{l r r r r r r}
\toprule
Variable & $N$ & Mean & Std & P25 & Median & P75 \\
\midrule
\multicolumn{7}{l}{\textit{Panel A: Returns}} \\[2pt]
$R_{i,w}$
  & 1{,}910{,}306 & $0.0048$ & $0.0813$ & $-0.0282$ & $0.0024$ & $0.0336$ \\
$R_{i,w+1}$
  & 1{,}910{,}306 & $0.0023$ & $0.0724$ & $-0.0290$ & $0.0016$ & $0.0321$ \\[5pt]
\multicolumn{7}{l}{\textit{Panel B: RSJ and decomposition}} \\[2pt]
RSJ
  & 1{,}910{,}306 & $0.0049$ & $0.1382$ & $-0.0845$ & $0.0047$ & $0.0942$ \\
$\hat\alpha_{\text{RSJ}}$ (M2 intercept)
  & 1{,}910{,}306 & $0.0012$ & $0.0236$ & $-0.0140$ & $0.0009$ & $0.0160$ \\
$\hat\beta_{\text{RSJ}}$ (M2 loading)
  & 1{,}910{,}306 & $0.3384$ & $0.2257$ & $0.1883$  & $0.3355$ & $0.4841$ \\[5pt]
\multicolumn{7}{l}{\textit{Panel C: Firm characteristics}} \\[2pt]
$\log(\text{ME})$
  & 1{,}910{,}306 & $14.119$ & $1.540$ & $12.990$ & $13.928$ & $15.048$ \\
BM
  & 1{,}910{,}306 & $0.569$  & $0.566$ & $0.265$  & $0.455$  & $0.725$  \\
MOM
  & 1{,}910{,}306 & $0.226$  & $0.847$ & $-0.136$ & $0.103$  & $0.380$  \\
REV
  & 1{,}910{,}306 & $0.0044$ & $0.0738$& $-0.0279$& $0.0023$ & $0.0330$ \\
IVOL
  & 1{,}910{,}306 & $0.0214$ & $0.0160$& $0.0114$ & $0.0173$ & $0.0266$ \\
ILLIQ (log)
  & 1{,}910{,}306 & $-20.099$& $1.975$ & $-21.463$& $-20.035$& $-18.676$\\[5pt]
\multicolumn{7}{l}{\textit{Panel D: Data quality}} \\[2pt]
Valid trading days / week
  & 1{,}910{,}306 & $4.844$ & $0.845$ & $5.000$ & $5.000$ & $5.000$ \\
Intraday obs.\ / stock-week
  & 1{,}910{,}306 & $4{,}587$ & $6{,}701$ & $879$ & $2{,}222$ & $5{,}516$ \\
\bottomrule
\end{tabular}
\begin{tablenotes}
\footnotesize
\item \textit{Notes:} Restricted to January 1993--December 2013 for comparability with
\citet{bollerslev2020}. $\hat\alpha_{\text{RSJ}}$ and $\hat\beta_{\text{RSJ}}$ are the
intercept and slope from a rolling 52-week regression of individual-stock RSJ on
market-wide RSJ (Method~2). Full sample in Table~\ref{tab:sumstats_full}.
\end{tablenotes}
\end{threeparttable}
\end{table}

\begin{table}[!htbp]
\centering
\caption{Supplementary Fama--MacBeth Specifications: Idiosyncratic Dominance (B9--B10)}
\label{tab:fm_supplementary}
\small
\begin{threeparttable}
\begin{tabular}{l c c}
\toprule
 & B9 & B10 \\
 & (RSJ idio $+$ RES idio) & (RSJ sys $+$ RES idio) \\
\midrule
RSJ$^{\text{idio}}$ (M1)
  & $-41.74^{***}$ & ---            \\
  & $(-4.27)$      &                \\[3pt]
RSJ$^{\text{sys}}$ (M1)
  & ---            & $-0.94$        \\
  &                & $(-0.22)$      \\[3pt]
RES$^{\text{idio}}$
  & $-54.71^{**}$  & $-54.71^{**}$  \\
  & $(-2.64)$      & $(-2.64)$      \\[5pt]
Controls (6)        & Yes & Yes \\
$N_{\text{weeks}}$  & 1{,}463 & 1{,}463 \\
\bottomrule
\end{tabular}
\begin{tablenotes}
\footnotesize
\item \textit{Notes:} Time-averaged Fama--MacBeth slope coefficients (bps/week) with
Newey--West $t$-statistics (six lags) in parentheses. B9 includes idiosyncratic RSJ
(M1) and idiosyncratic RES jointly with controls; B10 replaces idiosyncratic RSJ with
its systematic counterpart. Comparing B9 and B10 isolates the role of idiosyncratic
versus systematic RSJ conditional on idiosyncratic tail risk. Controls are identical to
those in Table~\ref{tab:fm_aggregate}. These specifications complement the primary full
decomposition results in Table~\ref{tab:fm_decomposition}. \\
$^{*}p<0.10$,\ $^{**}p<0.05$,\ $^{***}p<0.01$.
\end{tablenotes}
\end{threeparttable}
\end{table}

\begin{table}[!htbp]
\centering
\caption{Crisis-State Analysis: Complete Results Across All Specifications}
\label{tab:crisis_full}
\small
\begin{threeparttable}
\begin{tabular}{l l c c c c}
\toprule
Spec. & Predictor
  & Normal (bps) & Crisis (bps) & $\Delta$ (bps) & $t(\Delta)$ \\
\midrule
\multicolumn{6}{l}{\textit{Panel A: Univariate RSJ specifications}} \\[2pt]
RSJ (univ.) & RSJ Total & $-80.92^{***}$ & $-117.27^{**}$ & $-36.35$ & $-0.69$ \\
RSJ Sys M2 (univ.) & RSJ Sys (M2) & $133.02^{**}$ & $2.94$ & $-130.08$ & $-0.40$ \\
RSJ Idio M2 (univ.) & RSJ Idio (M2) & $-51.21^{***}$ & $-93.16^{*}$ & $-41.95$ & $-0.83$ \\
RSJ Sys M1 (univ.) & RSJ Sys (M1) & $-2.13$ & $20.20$ & $22.33$ & $0.78$ \\
RSJ Idio M1 (univ.) & RSJ Idio (M1) & $-43.04^{***}$ & $-93.70^{**}$ & $-50.66$ & $-1.05$ \\[5pt]
\multicolumn{6}{l}{\textit{Panel B: Univariate RES specifications}} \\[2pt]
RES (univ.) & RES Total & $-48.82^{*}$ & $98.55$ & $147.37$ & $1.41$ \\
RES Sys (univ.) & RES Systematic & $-180.85$ & $-158.23$ & $22.62$ & $0.04$ \\
RES Idio (univ.) & RES Idio & $-53.10^{**}$ & $102.61$ & $155.71$ & $1.61$ \\[5pt]
\multicolumn{6}{l}{\textit{Panel C: Controlled specifications (B1--B3)}} \\[2pt]
B1 $+$ ctrl & RSJ Total & $-74.55^{***}$ & $-92.82^{**}$ & $-18.27$ & $-0.40$ \\
B2 $+$ ctrl & RES Total & $-43.44^{*}$ & $55.22$ & $98.66$ & $1.31$ \\
B3 (joint) & RSJ Total & $-69.03^{***}$ & $-78.57^{*}$ & $-9.54$ & $-0.19$ \\
B3 (joint) & RES Total & $-63.38^{***}$ & $33.19$ & $96.57$ & $1.35$ \\[5pt]
\multicolumn{6}{l}{\textit{Panel D: Full decomposition --- M2 (B8)}} \\[2pt]
B8 & RSJ Idio (M2) & $-57.18^{***}$ & $-73.24$ & $-16.06$ & $-0.33$ \\
B8 & RSJ Sys (M2) & $20.22$ & $-10.43$ & $-30.65$ & $-0.19$ \\
B8 & RES Idio & $-78.49^{***}$ & $49.64$ & $128.13^{**}$ & $2.28$ \\
B8 & RES Sys & $-40.74$ & $-279.62$ & $-238.89$ & $-0.53$ \\[5pt]
\multicolumn{6}{l}{\textit{Panel E: Full decomposition --- M1 (B7)}} \\[2pt]
B7 & RSJ Idio (M1) & $-42.38^{***}$ & $-56.74^{*}$ & $-14.36$ & $-0.35$ \\
B7 & RSJ Sys (M1) & $-4.07$ & $5.97$ & $10.05$ & $0.42$ \\
B7 & RES Idio & $-66.94^{***}$ & $45.10$ & $112.04^{**}$ & $2.01$ \\
B7 & RES Sys & $-69.83$ & $-281.85$ & $-212.02$ & $-0.47$ \\
\bottomrule
\end{tabular}
\begin{tablenotes}
\footnotesize
\item \textit{Notes:} Normal and Crisis columns report the time-averaged FM slope
during non-recession and recession weeks, respectively. $\Delta$ is the
crisis-minus-normal difference; $t(\Delta)$ uses NW(6) standard errors.
NBER recessions: 2001Q1--Q4, 2007Q4--2009Q2, 2020Q1--Q2 (134 recession weeks).
$N_{\text{normal}}$: 1{,}277--1{,}481 depending on the specification. Controlled
specifications include the full control set from Table~\ref{tab:fm_aggregate}.
The primary crisis-state result is in Table~\ref{tab:crisis_state}.
$^{*}p<0.10$,\ $^{**}p<0.05$,\ $^{***}p<0.01$.
\end{tablenotes}
\end{threeparttable}
\end{table}

\begin{table}[!htbp]
\centering
\caption{Market Frictions: Complete Interaction Results (All Signal Variants)}
\label{tab:frictions_full}
\small
\begin{threeparttable}
\begin{tabular}{l c c c r}
\toprule
Signal
  & $\hat\beta_Z$ (bps) & $\hat\beta_{\text{ILLIQ}}$ (bps)
  & $\hat\beta_\times$ (bps) & $N_{\text{wks}}$ \\
\midrule
RSJ Total
  & $-444.55^{***}$ & $-7.56^{***}$ & $-18.73^{***}$ & 1{,}590 \\
  & $(-7.94)$ & $(-5.77)$ & $(-6.86)$ & \\[2pt]
RSJ Idiosyncratic (M2)
  & $-308.34^{***}$ & $-6.93^{***}$ & $-12.74^{***}$ & 1{,}411 \\
  & $(-6.09)$ & $(-4.91)$ & $(-5.08)$ & \\[2pt]
RSJ Systematic (M2)
  & $678.87^{***}$ & $-6.67^{***}$ & $31.27^{***}$ & 1{,}411 \\
  & $(2.89)$ & $(-4.72)$ & $(2.66)$ & \\[2pt]
RSJ Idiosyncratic (M1)
  & $-314.55^{***}$ & $-6.96^{***}$ & $-13.48^{***}$ & 1{,}463 \\
  & $(-4.00)$ & $(-4.83)$ & $(-3.65)$ & \\[2pt]
RSJ Systematic (M1)
  & $120.53^{***}$ & $-7.91^{***}$ & $6.41^{***}$ & 1{,}463 \\
  & $(3.17)$ & $(-5.12)$ & $(3.12)$ & \\[5pt]
RES Total
  & $-95.09$ & $-6.41^{***}$ & $-3.96$ & 1{,}462 \\
  & $(-0.74)$ & $(-3.57)$ & $(-0.56)$ & \\[2pt]
RES Idiosyncratic
  & $-82.24$ & $-6.58^{***}$ & $-3.07$ & 1{,}462 \\
  & $(-0.65)$ & $(-3.69)$ & $(-0.44)$ & \\[2pt]
RES Systematic
  & $-2225.43^{***}$ & $-5.49^{***}$ & $-107.06^{***}$ & 1{,}462 \\
  & $(-3.67)$ & $(-3.56)$ & $(-3.46)$ & \\
\bottomrule
\end{tabular}
\begin{tablenotes}
\footnotesize
\item \textit{Notes:} Each row reports a separate FM interaction regression including
the signal $Z$, ILLIQ, $Z\times\text{ILLIQ}$, and controls ($\log\text{ME}$, BM, MOM,
REV, IVOL). Coefficients in bps/week; $t$-statistics in parentheses use NW(6) standard
errors. The large RES Systematic coefficient reflects the small cross-sectional range of
that measure rather than a large economic effect. The key rows for the main-text
interpretation are reported in Table~\ref{tab:frictions}.
$^{*}p<0.10$,\ $^{**}p<0.05$,\ $^{***}p<0.01$.
\end{tablenotes}
\end{threeparttable}
\end{table}

\begin{table}[!htbp]
\centering
\caption{Signal-on-Friction Regressions: Idiosyncratic Characteristics and Illiquidity}
\label{tab:frictions_signal}
\small
\begin{threeparttable}
\begin{tabular}{l c c c c c}
\toprule
Signal
  & Univ.\ $\hat\delta$ & $t$ & Ctrl.\ $\hat\delta$ & $t$ & $N_{\text{wks}}$ \\
\midrule
RSJ Idiosyncratic (M2)
  & $0.0007^{***}$ & $(3.56)$ & $0.0027^{***}$ & $(6.96)$ & 1{,}412 \\[2pt]
RSJ Idiosyncratic (M1)
  & $-0.0005^{*}$ & $(-2.36)$ & $0.0042^{***}$ & $(9.22)$ & 1{,}464 \\[2pt]
RES Idiosyncratic
  & $0.0378^{***}$ & $(22.67)$ & $0.0163^{***}$ & $(28.92)$ & 1{,}463 \\
\bottomrule
\end{tabular}
\begin{tablenotes}
\footnotesize
\item \textit{Notes:} Weekly cross-sectional regressions of idiosyncratic signal on
log-Amihud illiquidity. The controlled specification adds $\log(\text{ME})$, BM, MOM,
REV, IVOL; ILLIQ is excluded from the control block. The sign reversal for RSJ
idiosyncratic (M1) --- negative univariate, positive controlled --- indicates that raw
negative correlation reflects offsetting size/volatility effects: controlling for firm
characteristics reveals positive friction-exposure comovement. RES idiosyncratic is
strongly and robustly positively associated with illiquidity in both specifications.
$t$-statistics use NW(6) standard errors.
$^{*}p<0.10$,\ $^{**}p<0.05$,\ $^{***}p<0.01$.
\end{tablenotes}
\end{threeparttable}
\end{table}

\clearpage
\bibliographystyle{apalike}
\bibliography{references}

\end{document}